% \documentclass[12pt]{article}
\documentclass[12pt, draft]{article}
% 
\usepackage{setspace}
\usepackage[labelfont=bf]{caption}
% \usepackage[margin=1in]{geometry}
% \usepackage[a4paper, total={6in, 10in}]{geometry}
\usepackage[margin=1in]{geometry}
\usepackage{amsmath,amssymb,amsthm,bm,mathtools}
\usepackage{algorithm}
\usepackage{algpseudocode}
\usepackage{enumitem}
\usepackage{hyperref}

% \documentclass{article}
\usepackage{graphicx} % Required for inserting images
\usepackage{hyperref}
% \usepackage[a4paper, total={5in, 10in}]{geometry}
% \usepackage{xcolor}
% \usepackage{color,soul}
\usepackage{listings}
\usepackage{booktabs}
% \usepackage{natbib}

% \usepackage{graphicx}
\usepackage{xcolor} % Required for colored text

% Define custom colors to match the image's muted green and red
\definecolor{impGreen}{rgb}{0.0, 0.5, 0.0}
\definecolor{worsRed}{rgb}{0.7, 0.0, 0.0}

% Citations
% \usepackage{natbib}
% Use the biblatex package, specifying the backend (biber is recommended)
\usepackage[style=numeric, sorting=nyt]{biblatex}
% Import your bibliography file (omit the .bib extension here)
\addbibresource{references.bib}


% \documentclass{article}
% \usepackage{booktabs}
% \usepackage{geometry}
% \geometry{a4paper, margin=1in}

% Math
\usepackage{amsmath}
\usepackage{amssymb}

% Subplots
\usepackage{graphicx}
\usepackage{subcaption}

% colored area
\usepackage{tcolorbox}
\usepackage{xcolor} % Required for defining custom colors
\definecolor{lightgray}{gray}{0.9} % Define a light gray color

% Highlights
\usepackage{color,soul}

% TODOs
% \usepackage{todonotes}
\setlength{\marginparwidth}{2cm}
\usepackage[textsize=tiny]{todonotes}
% Indicator
\usepackage{dsfont}

% Line numbers
\usepackage{lineno}
\renewcommand{\linenumberfont}{\normalfont\bfseries\tiny\color{gray}}
\linenumbers

%
\usepackage{multicol}

% Colors in table
% \usepackage{booktabs}
\usepackage{makecell}
\usepackage{xcolor}
% \usepackage{graphicx}

% 
% \documentclass{book} % Use the book or report class
\newtheorem{lemma}{Lemma}
\newtheorem{corollary}{Corollary}


\hypersetup{colorlinks=true,linkcolor=blue,citecolor=blue,urlcolor=blue}


\title{\vspace{-5mm}Fairness in Machine Learning for Real Estate Valuation:\\ Reducing Regressivity in Mass Appraisal \vspace{0mm}}
% Simple Regressivity Control for Mass Appraisal Pipelines
\author{Nicolás Acevedo Villena\\
{\small Operations Research Center, MIT, Cambridge, MA 02139. \texttt{nacevedo@mit.edu}}}
\date{
% \small
% Operations Research Center
\vspace{-8mm}
}

% theorem envs
\theoremstyle{plain}
\newtheorem{theorem}{Theorem}
\newtheorem{proposition}{Proposition}
\theoremstyle{definition}
\newtheorem{definition}{Definition}
\newtheorem{assumption}{Assumption}
\newtheorem{remark}{Remark}

\newcommand{\R}{\mathbb{R}}
\newcommand{\E}{\mathbb{E}}
\newcommand{\prox}{\mathrm{prox}}
\newcommand{\argmin}{\mathop{\mathrm{arg\,min}}}
\newcommand{\norm}[1]{\left\lVert #1 \right\rVert}
\newcommand{\ip}[2]{\left\langle #1,#2 \right\rangle}
\newcommand{\1}{\mathbf{1}}
\newcommand{\Cov}{\text{Cov}}
\newcommand{\Corr}{\text{Corr}}
\newcommand{\Var}{\text{Var}}


% Set commands
\newcommand{\Y}{\mathcal{Y}}
\newcommand{\X}{\mathcal{X}}
\newcommand{\V}{\mathcal{V}}
\newcommand{\U}{\mathcal{U}}
\newcommand{\G}{\mathcal{G}}
\newcommand{\D}{\mathcal{D}}
\newcommand{\N}{\mathcal{N}}
\newcommand{\F}{\mathcal{F}}
\newcommand{\A}{\mathcal{A}}
\newcommand{\J}{\mathcal{J}}
\newcommand{\I}{\mathcal{I}}
\newcommand{\M}{\mathcal{M}}
\newcommand{\C}{\mathcal{C}}
\renewcommand{\L}{\mathcal{L}}
\renewcommand{\H}{\mathcal{H}}
\renewcommand{\P}{\mathcal{P}}
\renewcommand{\R}{\mathbb{R}}

\newcommand{\diag}{\mathbf{\operatorname{diag}}}
% \newcommand{\R}{\mathbb{R}}
% \newcommand{\1}{\mathbf{1}}

\newcommand{\muBold}{\boldsymbol{\mu}}
\renewcommand{\t}{^{\top}}
\newcommand{\qt}[1]{``#1''}

\newcommand{\Expected}[1]{\mathbb{E}[#1]}
\newcommand{\Prob}[1]{\mathbb{P}(#1)}

\newcommand{\red}[1]{{\color{red} #1}}

% Math sets commands


\doublespacing
\begin{document}
\maketitle
% \tableofcontents

\begin{abstract}

Mass appraisal is a critical tool for local governments: it helps determine property taxes, one of the major sources of local public revenue in the United States. In recent years, assessors have increasingly adopted Machine Learning (ML) methods in the mass appraisal process. However, the ML-based property valuation literature has primarily focused on predictive accuracy, while giving less attention to equity concerns in assessment. One of the main equity concerns in mass appraisal is \textit{regressivity}, a form of vertical inequity, whereby lower-value homes tend to be over-assessed while higher-value homes tend to be under-assessed. We propose a simple assessor-compatible regressivity-aware modification of standard regression loss functions that can be incorporated into regression models commonly used in practice by assessors (e.g., linear regression \& gradient boosting). Rather than proposing a new valuation paradigm, the modification is designed as a targeted correction that can be integrated into existing assessor workflows with minimal disruption. Specifically, the modification consists of a covariance-based regularization term that softly penalizes the dependence between assessment ratios and property values. We empirically show that these soft-constrained models can substantially mitigate regressivity while incurring little to no loss in out-of-sample accuracy. Finally, we propose a multi-objective optimization framework for model selection that identifies Pareto-efficient tradeoffs among competing performance measures (e.g., accuracy \& regressivity) on out-of-sample data.

\end{abstract}

\section{Introduction}

Mass appraisal is a critical component of local public finance: it helps local governments determine property taxes, often the largest source of local public revenue in the United States \cite{Berry2021}. Mass appraisal refers to the process of estimating the value of a group of real estate properties, as of a given date, for tax purposes. The process is done using common data, standardized methods, and statistical testing \cite{IAAO2025Mass}. Its importance extends beyond property owners for two reasons. First, assessments help determine property tax liabilities, which in turn finance roads, schools, and many other public services.  Second, they affect property values and housing costs more broadly, with consequences that extend to much of the population, both directly and indirectly \cite{Berry2021, Rootetal2023}. Therefore, securing an \textit{accurate} and equitable (\textit{fair}) assessment process is essential for local public finance and social welfare.

In recent years, local government assessors' offices have increasingly adopted Machine Learning (ML)-based methods in mass appraisal, motivated by their strong predictive performance, efficiency, and scalability in large valuation settings \cite{IAAOAITaskForce2022,BidansetRakow2022,Rootetal2023}. Commonly studied model classes in the mass appraisal literature include linear regression-based methods and tree ensembles \cite{Rootetal2023,Ostrikova2024,Sevgen2024100111}. In this setting, however, predictive gains alone are not sufficient: methods must also remain explainable, operationally feasible, and compatible with assessor workflows. Naturally, the growing use of ML raises important equity concerns in assessment. Two of the main concerns are whether ML-based valuation can produce \textit{fair} predictions across: (i) properties with similar characteristics (\textit{horizontal equity}), and (ii) properties of different price level strata (\textit{vertical equity}) \cite{Quintos2014,10578d49-78f8-3ef9-8277-9b6add2bd071,cornia2006property}. Although both concepts have long been studied, even before the recent wave of ML adoption, vertical equity remains a persistent challenge in practice; assessments often exhibit \textit{regressivity} \cite{McMillen2020,IAAO2023VerticalReview,McMillenSingh2023}. Hence, this paper focuses primarily on reducing vertical equity problems, and in particular on reducing regressivity in ML-based valuation.

\textit{Regressivity} arises when lower-valued properties tend to be over-assessed, while higher-valued properties tend to be under-assessed. This pattern has long been recognized in the mass appraisal literature and can have substantial distributional consequences for taxpayers \cite{McMillen2020}. \textcite{amornsiripanitch2020why} argues that owners of
lower-valued properties might face effective tax rates up to 50\% higher than those paid by owners of higher-valued properties. Data limitations, model misspecification, and assessment practices are often cited as contributing factors \cite{Berry2021}. The adoption of ML-based assessment workflows creates both an opportunity and a challenge in this setting. On one hand, these workflows provide a mathematical object that can be explicitly modified to mitigate regressivity (see, e.g., \textcite{CandoganHanLu2023}). On the other hand, standard objectives in regression models are designed to optimize average predictive accuracy, which tends to pull extreme predictions toward conditional averages. Thus, if valuation features are incomplete or imperfectly measured in the mass appraisal context, predictive models can reinforce regressive assessment patterns \cite{amornsiripanitch2020why, CandoganHanLu2023}.

% by overvaluing lower-priced properties and undervaluing higher-priced ones 

To evaluate assessment equity, we follow the standard framework used in ratio studies by the International Association of Assessing Officers (IAAO) \cite{IAAO2025ExposureRatio}. Many of these measures are based on the \textit{assessment-to-sale ratio}, defined as $r := \hat{P}/P$, where $\hat{P}$ denotes the assessed value and $P$ the market value of a property. Under regressive assessment, these ratios tend to decrease with property value: lower-tier properties are more likely to have ratios above one, while higher-tier properties are more likely to have ratios below one. Figure~\ref{fig:regressivity_motivation} illustrates this phenomenon. Motivated by this pattern, we study an in-processing approach that penalizes systematic dependence between assessment ratios and property values, with the goal of reducing regressivity while preserving predictive accuracy.

\begin{figure}[!h]
    \centering 
    \includegraphics[width=0.7\linewidth]{Fairness Price/CCAO meeting/img/Baseline LightGBM_0_motivation.pdf}
    \caption{Illustration of regressive behavior in mass appraisal. Assessment-to-sale ratios $r$ tend to decrease with property value $P$: lower-value properties are more likely to be over-assessed ($r > 1$), while higher-value properties are more likely to be under-assessed ($r < 1$).}
    \label{fig:regressivity_motivation}
\end{figure}

More broadly, this paper helps bridge two bodies of work that have largely developed in parallel: the ML-based mass-appraisal literature, which has mainly emphasized predictive performance model comparison in property valuation \cite{Rootetal2023,Zilli2024}, and the fairness-in-ML literature, which has developed general frameworks for fair regression, fairness regularization, and fairness--accuracy tradeoff analysis in predictive models \cite{CatonHaas2024,AgarwalDudikWu2019,FitzsimonsEtAl2019,lee2022maximal,LiEtAl2022Fairness}. This connection is especially natural in mass appraisal because regressivity can be viewed as a systematic dependence between assessment ratios and property values, allowing an assessor-defined notion of vertical inequity \cite{IAAO2025ExposureRatio} to be addressed using tools from fairness-aware regression \cite{CandoganHanLu2023}.

\noindent \textbf{Contributions.} Our main contributions are as follows:
\begin{itemize}
    \item We develop a simple, in-processing approach for mitigating regressivity in ML-based mass appraisal by augmenting standard regression losses with a soft covariance-based regularization term involving transformations of the assessment-to-sales ratio $r$ and the property market value $P$.\footnote{Related dependence-regularization ideas appear in the fair-ML literature; see, e.g., \cite{lee2022maximal,LiEtAl2022Fairness,ScutariPaneroProissl2022}.}
    
    \item We show how this approach can be incorporated into regression models commonly used in practice by assessors, including linear regression and gradient boosting, thereby providing a practical way to reduce vertical inequity without requiring a complete redesign of existing valuation workflows.
    
    \item Using a real-world mass-appraisal pipeline from the Cook County Assessor's Office, we show empirically that the proposed corrections can improve equity--accuracy tradeoffs under standard assessment metrics, and in some settings improve both dimensions simultaneously.
    
    \item We formalize a multi-objective model-selection framework for calibration that targets Pareto-efficient tradeoffs between predictive accuracy and equity measures. In particular, the framework can recover solutions that prioritize predictive performance subject to equity requirements, and it also admits an extension based on convex combinations of candidate models to achieve improved tradeoffs. \footnote{For related Pareto/fairness-frontier formulations and model-stacking ideas, see, e.g., \cite{martinez2020minimax,little2022fairness}.}
\end{itemize}

The remainder of the paper reviews the relevant background on mass-appraisal equity and fairness-aware regression, develops the proposed regularization and model-selection framework, and evaluates the resulting methods on a real-world assessment pipeline.




\section{Related Work}

\subsection{ML in Mass Appraisal and Assessor Workflows}

Recent reviews show that mass appraisal and real-estate valuation research have moved steadily toward automated valuation models (AVMs) and machine learning (ML) methods, with multiple regression, tree ensembles, and neural-network-based models among the most frequently studied approaches \cite{WangLi2019,Rootetal2023,Zilli2024,purdueAssessingTechniques}. Across this literature, model evaluation is still dominated by predictive criteria---typically accuracy and model comparison across algorithms---rather than by equity as an explicit model-design target \cite{WangLi2019,Rootetal2023,Zilli2024}. In particular, recent reviews identify Random Forest, Gradient Boosting, and related ensemble methods as some of the most commonly used and best-performing model classes in modern mass appraisal studies \cite{Rootetal2023,Sevgen2024100111,Zilli2024}. This accuracy-centered orientation leaves open a central question for the present paper: how should modern valuation models be designed and selected when equity is also a primary objective?

From an assessor perspective, however, adoption depends on more than predictive fit alone. The assessor-oriented AI literature emphasizes explainability, defensibility, administrative feasibility, and integration into office workflows as central conditions for real-world use \cite{IAAOAITaskForce2022}. Consistent with this, \textcite{BidansetRakow2022} report that AVM adoption in assessment offices remains partial and uneven, with perceived barriers including communication challenges, technical capacity, and budget constraints. As a result, the relevant literature for the present paper is not only the literature on which models predict most accurately, but also the literature on how advanced models can be incorporated into assessor practice in a way that remains operationally useful and publicly defensible \cite{IAAOAITaskForce2022,BidansetRakow2022}. These practical constraints make low-disruption mitigation strategies especially relevant in mass appraisal: an equity intervention is more likely to matter in practice when it preserves familiar model classes and fits within existing office workflows.





\subsection{Equity in Mass Appraisal: Horizontal Equity, Vertical Equity, and Regressivity}

Equity has long been treated as a core objective of mass appraisal rather than as a secondary byproduct of predictive performance \cite{IAAO2013Ratio,IAAO2025Mass,Quintos2014}. In assessor practice, this objective is commonly decomposed into \emph{horizontal equity}, which concerns the uniform treatment of similarly situated properties, and \emph{vertical equity}, which concerns the absence of systematic price-related bias across the value distribution \cite{10578d49-78f8-3ef9-8277-9b6add2bd071,cornia2006property,McMillenSingh2023}. The basic diagnostic framework remains the sales-ratio study: current IAAO standards continue to emphasize ratio-study analysis, stratified evaluation across relevant submarkets or property classes, and careful sales verification as the foundation for evaluating mass-appraisal performance \cite{IAAO2013Ratio,IAAO2020SalesVerification,IAAO2025Mass}. Within this framework, the coefficient of dispersion (COD) remains a standard measure of uniformity, while the price-related differential (PRD) and price-related bias (PRB) remain central diagnostics for vertical equity \cite{IAAO2013Ratio,IAAO2025Mass}.

At the same time, the recent assessor literature has become more cautious about relying on any single measure of vertical equity. \textcite{IAAO2023VerticalReview} compare the behavior of several vertical-equity diagnostics and recommend reporting a \emph{suite} of tests rather than privileging one statistic. \textcite{McMillenSingh2023} further show that some regression-based procedures for detecting regressivity can be severely biased when assessments themselves are generated from regression-based valuation models, and they argue for complementing classical statistics such as the PRD with broader distribution-based diagnostics. This concern is already influencing standards discussions: the September 2025 IAAO exposure draft introduces the Vertical Equity Indicator (VEI) as an emerging principal test while still emphasizing that no single diagnostic can detect all forms of vertical inequity \cite{IAAO2025ExposureRatio}.

Beyond diagnostics, recent policy-oriented work shows that regressivity is not an isolated local anomaly but a widespread and substantively important feature of property taxation \cite{Berry2021,AvenancioLeonHoward2022}. Using millions of residential sales, \textcite{Berry2021} finds that assessments are typically regressive, that this pattern is widespread across U.S.\ counties, and that it is more consistent with limitations in assessors' data and models than with simple measurement error or explicit policy choices alone. \textcite{AvenancioLeonHoward2022} broaden the discussion by documenting a nationwide assessment gap and showing that assessment inequities can translate into systematic differences in tax burdens for the same bundle of local public services. Assessor-oriented studies also emphasize that institutional design and empirical design matter: \textcite{ScottDaigle2022} document a program that reduced both vertical and horizontal inequity, while \textcite{Rakow2023} and \textcite{Rakow2024} show that measured local regressivity and its effective incidence can depend on assessment timing, exemptions, and policy design.

A smaller but especially relevant strand of the assessor literature moves from equity measurement to equity correction within the valuation pipeline itself. \textcite{Quintos2014} propose a two-step spatial-lag/spatial-quantile procedure that improves both horizontal and vertical equity in a heterogeneous housing market. \textcite{BidansetEtAl2019} show that geographically weighted regression can reveal local vertical and horizontal inequities that are obscured by jurisdiction-wide conditional-mean models, and \textcite{BidansetEtAl2025} extend this line by using geographically weighted regression and clustering to construct more coherent market areas for mass valuation. Related work has also shifted attention from model fitting to model governance: \textcite{HermansEtAl2022} incorporate property-level ratio studies into AVM validation, and \textcite{HermansEtAl2025} use Monte Carlo sampling to assess the probability that an AVM will fail ratio-study standards under alternative sampling configurations. At the research frontier, \textcite{CandoganHanLu2023} explicitly study the accuracy--fairness trade-off in regressive property taxation using segmented, fairness-aware valuation models. Taken together, this literature shows that equity in mass appraisal does not arise automatically from better prediction alone; rather, it increasingly requires explicit measurement, explicit monitoring, and, in some settings, explicit correction within the modeling and calibration process. For the present paper, the key implication is that vertical inequity is already understood by assessors as a systematic pattern in the relationship between assessment outcomes and value level, which makes it a natural target for methods that explicitly control such relationships during training.

\subsection{Fairness in ML for Regression}

Most of the fairness-in-ML literature was originally developed for classification problems, but many real applications involve continuous prediction rather than binary decisions \cite{CatonHaas2024,BarocasHardtNarayanan2023}. In those settings, fairness cannot usually be reduced to a single thresholded decision without discarding information about score magnitudes, residual structure, and calibration. For that reason, a substantial literature has developed around \emph{fair regression}: methods that treat fairness as an explicit objective when the prediction target is continuous \cite{AgarwalDudikWu2019,FitzsimonsEtAl2019,ChzhenEtAl2020Wasserstein,ChzhenEtAl2020Recalibration,WeiEtAl2023}.

A recurring conclusion of this literature is that no single fairness notion is universally preferred for regression. Depending on the application, one may want parity in the marginal distribution of predictions, parity of losses or residuals across groups, equality of conditional means, or some broader structural notion of non-discrimination \cite{AgarwalDudikWu2019,CatonHaas2024,BarocasHardtNarayanan2023}. Methodologically, the state of the art has moved from ad hoc heuristics toward constrained or regularized empirical risk minimization. \textcite{AgarwalDudikWu2019} formulate fair regression under statistical-parity and bounded-group-loss constraints and provide reduction-based algorithms with guarantees for broad classes of losses. \textcite{FitzsimonsEtAl2019} introduce a general framework based on group fairness in expectation that can be integrated into a variety of regression learners. \textcite{WeiEtAl2023} develop an RKHS formulation for mean-parity fair regression with multiple sensitive attributes. A complementary line studies distributional fairness through optimal transport and post-processing: \textcite{ChzhenEtAl2020Wasserstein} characterize demographic-parity regression via Wasserstein barycenters, while \textcite{ChzhenEtAl2020Recalibration} derive a plug-in-and-recalibration approach with statistical guarantees. Regression-specific optimization work also includes explicit dependence constraints: \textcite{KomiyamaEtAl2018} measure unfairness through correlation or coefficient-of-determination constraints, and \textcite{ScutariPaneroProissl2022} show that fairness can be controlled through a simple ridge-type penalization strategy. Taken together, this literature establishes a useful perspective for the present paper: in continuous prediction problems, fairness can be built directly into the training or calibration process rather than being treated only as an ex-post diagnostic.

\subsection{Dependence-Based Fairness Regularization and Fair Regression with Continuous Attributes}

A particularly relevant idea from fair regression is that unfairness can often be represented as an undesirable statistical relationship between model outputs and some variable of concern; mitigation then proceeds by explicitly penalizing that relationship during training. For the purposes of the present paper, the most relevant strand is therefore the one that interprets fairness as a problem of controlling statistical dependence between predictions and sensitive variables---or, more generally, between predictions and variables along which undesirable dependence is to be controlled. This perspective is especially important when the relevant source of unfairness is not naturally binary or categorical.

Several lines of work develop this idea in different ways. \textcite{MaryEtAl2019} extend fairness-aware learning to continuous sensitive attributes and treatments using the R\'enyi maximal correlation coefficient as a differentiable dependence measure. \textcite{KomiyamaEtAl2018} formulate regression fairness directly through correlation and coefficient-of-determination constraints. \textcite{lee2022maximal} further develop maximal-correlation-based regularization as a general framework for imposing fairness criteria through dependence control. In parallel, \textcite{PerezSuayEtAl2017} and \textcite{LiEtAl2022Fairness} use kernel-based dependence measures, especially HSIC and related cross-covariance operators, to build tractable dependence-regularized learners, and \textcite{BuylDefranceDeBie2024FAIRRET} provide a recent modular framework for differentiable fairness regularization terms.

This literature is especially relevant here for two reasons. First, dependence-based regularizers are often attractive in practice because they can be incorporated into standard optimization pipelines with relatively limited changes to the underlying learner \cite{KomiyamaEtAl2018,MaryEtAl2019,LiEtAl2022Fairness,lee2022maximal,BuylDefranceDeBie2024FAIRRET}. Second, they naturally accommodate settings in which the relevant source of inequity is encoded through a continuous or structurally important variable rather than a classical protected demographic attribute. In mass appraisal, the relevant dependence is not primarily between predictions and a demographic group indicator, but between assessment behavior and value level. This makes dependence-based fair-regression methods methodologically closer to the present problem than most of the better-known group-fairness methods designed for discrete protected classes. In that sense, fair-regression theory offers assessors not just a new vocabulary, but also a practical toolkit for turning an observed inequity pattern into an explicit modeling target.

\subsection{Fairness--Accuracy Tradeoffs and Multi-Objective Model Selection}

A separate but closely related literature treats fairness mitigation as a genuinely multi-objective problem rather than as a single constrained-learning problem with one fixed penalty parameter. At a high level, this view emphasizes that different models, penalties, and post-processing choices trace out different combinations of predictive performance and fairness, so model selection should often be understood as choosing among tradeoffs rather than identifying a single universally best model \cite{CatonHaas2024}. \textcite{martinez2020minimax} formulate fairness explicitly as a multi-objective optimization problem and characterize Pareto-efficient solutions that avoid unnecessary harm across groups. \textcite{little2022fairness} study the empirical fairness--accuracy Pareto frontier directly, proposing tools to identify, quantify, and optimize model tradeoffs through frontier-based summaries rather than one-number comparisons.

This perspective is especially natural for the present application. In assessor workflows, model calibration already requires comparing many candidate specifications across multiple out-of-sample criteria, such as accuracy, dispersion, and jurisdiction-specific diagnostics. Fairness-aware mass appraisal does not create a wholly new decision problem; rather, it extends an existing calibration problem by adding equity as an additional objective. The relevant methodological question is therefore not only how to construct a regressivity-aware learner, but also how to calibrate and select candidate models when accuracy and equity move in different directions. The Pareto and frontier literature provides the natural OR language for this problem and clarifies why model selection should itself be treated as part of the fairness-aware workflow.

\subsection{Positioning of the Present Work}

The present paper sits at the intersection of these two bodies of work. On the one hand, the mass-appraisal literature has developed rich standards, diagnostics, and application-specific corrections for horizontal and vertical equity, but it still treats regressivity primarily as a ratio-study outcome or as a problem to be corrected through specific spatial, segmented, or post hoc procedures \cite{Quintos2014,IAAO2023VerticalReview,BidansetEtAl2019,CandoganHanLu2023}. On the other hand, the fair-regression literature has developed general constrained and regularized approaches for continuous prediction, including dependence-based penalties and fairness--accuracy frontier methods, but typically in settings defined by demographic protected attributes rather than assessor-style notions of vertical inequity \cite{AgarwalDudikWu2019,LiEtAl2022Fairness,lee2022maximal,martinez2020minimax,little2022fairness}.

Our contribution is to connect these two literatures in an assessor-compatible way. We start from an assessor-defined problem---regressivity as a systematic pattern in the relationship between assessment ratios and value level---and recast it in a form that fair-regression tools can act on directly during training. This shifts the role of vertical-equity analysis from purely ex-post diagnosis toward in-processing control. Concretely, we use this interpretation to build a regularization approach for standard regression models used in mass appraisal and embed the resulting corrections in a multi-objective calibration framework based on accuracy--equity tradeoffs. In this sense, the paper does not introduce equity concerns into mass appraisal, nor does it introduce fairness regularization into machine learning; rather, it develops a practical bridge between assessor-defined vertical equity and fairness-aware regression methods.


\section{Preliminaries: Problem Formulation, Equity Measures, and Temporal Evaluation}
\label{sec:problem_setting}

This section fixes the basic prediction notation, defines the assessor-facing equity objects used throughout the paper, and formalizes the temporal evaluation setting in which models are trained, calibrated, and compared. The goal is to describe the predictive problem in the same language used by both modern regression workflows and ratio-study practice, so that the methodological section can build on a single, stable set of definitions.

\subsection{Mass-Appraisal Prediction Setting}
\label{subsec:prediction_setting}

Mass appraisal estimates the value of a \emph{group} of properties as of a given date using common data, standardized methods, and statistical testing \cite{IAAO2025Mass}. In the predictive setting studied here, we observe a sample of $n$ sold properties
$
\mathcal{D}_n := \{(x_i,P_i,t_i)\}_{i=1}^n,
$
where $x_i \in \mathcal{X}$ is the vector of observed characteristics for property $i \in [n]$ (e.g., squared footage, number of rooms, distance to a bus stop, etc.), $P_i \in \mathbb{R}_{++}$ is its verified sale price, and $t_i$ is its sale date. Because prices are positive and strongly right-skewed, we work on the log-price scale. That is, for each sold property $i$, define its log-price
$
y_i := \log P_i.
$

Let $\H$ denote some hypothesis class for a family of parametric regression models, and let
$
f:\mathcal{X}\to\mathbb{R},
$ with $ f \in \H,
$
denote a candidate predictor of log-price (e.g., linear regression or boosted trees). For a given property $i\in [n]$, the model produces the log-scale prediction
$
\hat y_i := f(x_i),
$
and the corresponding prediction on the original price scale is
$
\widehat P_i := \exp(\hat y_i).
$ 

The \textit{baseline} learning problem is the regularized empirical risk minimization problem
\begin{equation}
\label{eq:baseline_erm}
f_\mathrm{base}
\in
\arg\min_{f\in\H}
\left\{
\frac{1}{n}\sum_{i=1}^n \ell\big(f(x_i),y_i\big)
+
\Omega
\right\}
,
\end{equation}
where $\ell$ is a loss function on the log-price scale (e.g., squared error) and $\Omega$ is an optional complexity penalty or model-specific regularizer (e.g., $\ell_1$-norm for the LASSO). This formulation covers the main model classes studied later, including regularized linear regression and boosting-based learners.

\paragraph{Prediction scale and retransformation.}
All assessor-facing diagnostics in the paper are computed on the original sale price scale, so the predictive object ultimately evaluated is \(\widehat P_i=\exp(\hat y_i)\). This choice is deliberate: the target in the present workflow is the price-scale assessment ratio, not a bias-corrected estimate of \(\E[P\mid X]\), for example, under a lognormal model\todo{Add a citation as an example maybe (?)}.  The method is trained on the log-scale and then evaluated on the price scale, where ratio studies and related diagnosis are defined (see section \ref{subsec:equity_metrics}). If the reader prefers another transformation, the same framework can be applied under other mappings, but the central dependence-control logic introduced later is unchanged.


\subsection{Assessment Ratios and Equity Measures}
\label{subsec:equity_metrics}



For any fitted model $f\in \H$ and any evaluation sample $S$, define the \emph{assessment-to-sale ratio}
\begin{equation}
\label{eq:ratio_def}
r_i := \frac{\widehat P_i}{P_i} ,
\qquad \forall i\in S,
\end{equation}
where dependence on the fitted model is suppressed when clear from context. 

Ratio studies are conducted on the original price scale, as defined by the IAAO \cite{IAAO2025ExposureRatio}, so all equity diagnostics defined below are computed using $\hat P_i$ and $P_i$, and not directly the log-scale values $\hat y_i$ or $y_i$. Moreover, all of the metrics defined depend on the assessment ratio $r_i$, and the basic statistics around it. 

\paragraph{Assessment Level.} To measure the \textit{assessment level} \cite{IAAO2025ExposureRatio}, assessors look at the central statistics of the ratios; the median, average, weighted average, and standard deviation of $\{r_i\}_{i\in S}$.
Let
$
m_S := \operatorname{med}\{r_i: i\in S\}
$
denote the median ratio, let $\Bar{r}:= \frac{1}{|S|}\sum_{i\in S} r_i$ denote the average ratio, let $\overline{r_W}:=\frac{1}{\sum_{i\in S}P_i}\sum_{i\in S}r_i P_i$ be the value-weighted average ratio, and let $\sigma_S:=\sqrt{\frac{1}{|S|-1}\sum_{i\in S} (r_i - \bar r)^2}$ be the standard deviation of  the sample ratios $\{r_i\}_{i \in S}$.

\medskip


\noindent \textbf{Horizontal equity measures.} We summarize \textit{horizontal equity} primarily through the coefficient of dispersion (COD), which the IAAO describes as the average percentage deviation from the median ratio $m_S$ \cite{IAAO2013Ratio,IAAO2025ExposureRatio}:
\begin{equation}
\label{eq:cod}
\operatorname{COD}(S)
:=
100\cdot\frac{1}{|S|}
\sum_{i\in S} \frac{|r_i-m_S|}{m_S}.
\end{equation}
Lower COD values indicate greater within-group uniformity.  We also report coefficient of variation (COV), defined as \begin{equation}    
\mathrm{COV}(S):=100 \cdot \frac{\sigma_S}{\bar r},
\end{equation}
as a supplementary dispersion diagnostic. Both indicators measure deviation of the ratios $\{r_i\}_{i \in S}$ from one of their central statistics in $S$. 

% With an ideal predictor, there would be no deviation, and  would tend to 0\%.

\medskip

\noindent \textbf{Vertical equity measures.} We summarize vertical equity using \textit{price-related diagnostics}. First, the price-related differential (PRD) defined as
\begin{equation}
\label{eq:prd}
\operatorname{PRD}(S)
:=
\frac{\bar r}{\overline{r_W}}\;,
\end{equation}
that is, the mean ratio divided by the value-weighted mean ratio. Under standard assessor interpretation, PRD values above $1.03$ tend to indicate regressivity, while values below $0.98$ tend to indicate progressivity \cite{IAAO2013Ratio,IAAO2025ExposureRatio}.

We also report the coefficient of price-related bias (PRB), which depends on a quantity refer to as the \textit{market value proxy}; defined as $\operatorname{MV}_i := \log_2\left(\frac{1}{2} (P_i + \widehat{P}_i / m_S)\right)$, designed to reduce measurement-error bias, according to the IAAO specifications \cite{IAAO2025ExposureRatio}. The PRB is defined then as 
% \vspace{-2mm}
\begin{equation}
\operatorname{PRB}(S)=\text{slope}\left(\frac{r_i- m_S}{m_S}\sim
  \operatorname{MV}_i\right).
\end{equation}
That is, the slope of relative ratio deviations from the median against the market value proxy. PRB values below -0.05 suggest {regressivity}, while values above 0.05 suggest {progressivity} \cite{IAAO2013Ratio,IAAO2025ExposureRatio}.

Finally, we define the latest introduces measure by the IAAO, the vertical equity indicator (VEI), defined as 
\begin{equation}
\operatorname{VEI}(S)=100\cdot \left(\frac{m_{S{q_\text{10}}} - m_{S{q_\text{1}}}}{m_S}\right),
\end{equation}
where $m_{S{q_\text{10}}}$ is the median of the 10th decile of samples sorted by the proxy MV on $S$ (10\% highest values), and $m_{S{q_\text{1}}}$ is the median of the respective first decile (10\% lowest values). VEI values below -10\% suggest regressivity, while values above 10\% suggest progressivity \cite{IAAO2025ExposureRatio}.

Table \ref{tab:assessment_metrics_summary} summarizes the relevant metrics used in by assessors, and we use them in the empirical analysis and comparison of models, to measure the equity and beahavior of the predictions in Section \todo{reference}.



\begin{table}[!ht]
\centering
\caption{Summary of core assessment ratio statistics, desired targets, and interpretation.}
\label{tab:assessment_metrics_summary}
\begin{tabular}{lcp{8.5cm}}
\toprule
\textbf{Measure} & \textbf{Target Value} & \textbf{Interpretation} \\
\midrule
\multicolumn{3}{l}{\textit{{Central Statistics (Assessment Level)}}} \vspace{3mm} \\
Median ($m_S$) & \(\approx 1\) & Outlier-robust measure of overall assessment level. \\
Mean ($\Bar r$) & \(\approx 1\) & Arithmetic average; highly sensitive to extreme outliers. \\
Weighted Mean $(\overline{r_W})$ & \(\approx 1\) & Value-weighted average; serves as the denominator for calculating the PRD. \\
\midrule
\multicolumn{3}{l}{\textit{{Horizontal Equity (Uniformity)}}} \vspace{3mm}\\
COD & \([5.0, 15.0]\textsuperscript{*}\) & Average \% deviation from the median. Lower values indicate better uniformity. \\
COV & $-$ & Ratio of standard deviation to the mean. Lower values mean less relative variability. \\
\midrule
\multicolumn{3}{l}{\textit{{Vertical Equity (Regressivity)}}} \vspace{3mm}\\
PRD & \([0.98, 1.03]\) &Values \(>1.03\) indicate \textit{regressivity} (low-value properties are relatively overassessed). \\
PRB & \([-0.05, 0.05]\) & Values \(<-0.05\) indicate \textit{regressivity} (\% bias per doubling of property value). \\
VEI & \([-10\%, 10\%]\) & Values \(<0\) indicate \textit{regressivity} (negative gap between the top and bottom value deciles). \\
\bottomrule
\multicolumn{3}{l}{\footnotesize \textsuperscript{*}Assessments with COD values below 5.0 are considered highly suspect. Acceptable COD upper limits} \\
\multicolumn{3}{l}{\footnotesize  {\color{white}\textsuperscript{*}}vary by property type (up to 20.0 or 25.0 for rural/vacant) \cite{IAAO2025ExposureRatio}. } \\
\end{tabular}
\end{table}

% \begin{table}[!ht]
% \centering
% \caption{Reference interpretation ranges used in the empirical discussion.}
% \label{tab:reference_ranges}
% \begin{tabular}{lcc}
% \toprule
% Measure & Preferred region & Interpretation \\
% \midrule 
% PRD & \([0.98,1.03]\) & Values \(>1.03\) suggest regressivity \\
% PRB & \([-0.05,0.05]\) & Closer to \(0\) indicates less price-related bias \\
% VEI & \([-10\%,10\%]\) & Closer to \(0\) indicates less value-tier distortion \\
% % MKI & \(\approx 1\) & Closer to \(1\) indicates more proportional assessment \\
% \bottomrule
% \end{tabular}
% \end{table}


% The key point for this section is therefore simple: assessor-side equity is evaluated through ratio-based summaries on the original price scale, while the learning problem is solved on the log-price scale. The identity in \eqref{eq:log_ratio_identity} is what allows these two views to be connected.


{
\color{gray}
\subsection{Regressivity as a Dependence Problem}
\label{subsec:dependence_problem}

The assessor literature usually describes regressivity as a \emph{price-related} pattern in the ratios: as property value increases, assessment ratios tend to fall. We view this as a dependence problem between the assessment ratios and property value. This behavior can be identified in Figure \ref{fig:regressivity_motivation_boxed} from the negative trend between the ratios $\{r_i\}_{i\in S}$ and the log-prices $\{y_i\}_{i \in S}$. We use the log-scale in the figure because property values are typically highly skewed, so plotting against log-price makes the price-related trend more visually interpretable across the market.

\begin{figure}[!h]
    \centering 
    \includegraphics[width=0.7\linewidth]{Fairness Price/CCAO meeting/img/motivation/Baseline LightGBM_0_motivation_squared_only.pdf}
    \caption{Illustration of regressive behavior in mass appraisal. Assessment-to-sale ratios $r$ tend to decrease with property value $P$: lower-value properties are more likely to be over-assessed ($r > 1$), while higher-value properties are more likely to be under-assessed ($r < 1$).}
    \label{fig:regressivity_motivation_boxed}
\end{figure}

Since both the ratios and prices are continuous quantities, a natural way to measure their linear dependence is the covariance, defined empirically as
\begin{equation}
\label{eq:sample_cov}
\operatorname{Cov}(r,P)
:=
\frac{1}{|S|}
\sum_{i\in S}
\big(r_i-\bar r\big)\big(P_i-\bar P\big),
\end{equation}
where $\bar r$ and $\bar P$ denote the sample means of $\{r_i\}_{i\in S}$ and $\{P_i\}_{i\in S}$, respectively. In particular, when $\Var(r)>0$ and $\Var(P)>0$, the sign of the covariance determines the sign of both the correlation and the slope from the simple linear regression of $r$ on $P$:
\begin{equation}
\label{eq:cov_corr_slope_sign}
\operatorname{sign}\!\left(\Cov(r,P)\right)
=
\operatorname{sign}\!\left(\Corr(r,P)\right)
=
\operatorname{sign}\!\left(\beta_{r\mid P}\right),
\end{equation}
where $\operatorname{sign}(x)\in \{-1,0,1\}$ is the signum function, and correlation and slope are simply
\[
\Corr(r,P):=\frac{\Cov(r,P)}{\sqrt{\Var(r)\Var(P)}},
\qquad
\beta_{r\mid P}:=\frac{\Cov(r,P)}{\Var(P)}.
\]
Thus, $\Cov(r,P)<0$ means that ratios tend to decrease as prices increase, which is precisely the linear price-related pattern associated with regressivity.

Moreover, covariance is aligned with classical assessor-side regressivity diagnostics. In particular, it has the same directional ordering as the price-related differential (PRD) \cite{IAAO2025ExposureRatio}, introduced in Section \ref{subsec:equity_metrics}.

\begin{proposition}
\label{prop:cov_prd}
Let $P>0$ and $r=\widehat P/P>0$ almost surely, with finite first moments. Then the population covariance between $r$ and $P$, $\Cov(r,P)=\Expected{rP}-\Expected{r}\Expected{P},$ and the population $
\mathrm{PRD}=\frac{\Expected{r}\Expected{P}}{\Expected{rP}}$
satisfy the directional ordering
\begin{equation}
\label{eq:cov_prd_sign}
\operatorname{sign}\!\left(\Cov(r,P)\right)
=
\operatorname{sign}\!\left(1-\mathrm{PRD}\right).
\end{equation}
\end{proposition}

\begin{proof}
Since
$
\Cov(r,P)=\Expected{rP}-\Expected{r}\Expected{P},
$
we have that
\[
1-\mathrm{PRD}
=
1-\frac{\Expected{r}\Expected{P}}{\Expected{rP}}
=
\frac{\Expected{rP}-\Expected{r}\Expected{P}}{\Expected{rP}}
=
\frac{\Cov(r,P)}{\Expected{rP}}.
\]
Because $r>0$ and $P>0$ almost surely, it follows that $\Expected{rP}>0$. Therefore, $1-\mathrm{PRD}$ and $\Cov(r,P)$ have the same sign.
\end{proof}

% \noindent This proposition shows that covariance is not just an ad-hoc arbitrary dependence measure, but one that is directionally consistent with a standard price-related diagnostic used in assessment practice. For this reason, covariance provides a natural starting point for regressivity correction.

\noindent Proposition \ref{prop:cov_prd} shows that covariance is not just an arbitrary ad-hoc dependence measure, but one that is directionally consistent with a standard price-related diagnostic used in assessment practice. This makes $\Cov(r,P)$ a natural conceptual starting point for regressivity correction.

\noindent \textbf{Log-space residuals and covariance.} In later sections, however, we work on the log-price scale for model estimation. Hence, we introduce the definitions in detail for the log-scale. Recall the log-price
$
y:=\log P$ and the log-price predictions $\hat y:=\log \widehat P$ definitions. Let $e:=\hat y-y$ be the \textit{log-scale residuals}. 
A key identity is that the log-scale residuals $\{e_i\}_{i\in S}$ are exactly the logarithm of the assessment ratios $\{r_i\}_{i \in S}$. That is \vspace{-2mm}
\begin{equation}
\label{eq:log_ratio_identity}
e_i
=
\hat y_i-y_i
=
\log \widehat P_i-\log P_i
=
\log\!\left(\frac{\widehat P_i}{P_i}\right)
=
\log r_i,
\qquad \forall i\in S.
\end{equation}
Thus, the residuals in log-price space are not merely a proxy for the assessment ratios; it provides an exact log-scale representation of them. In particular,
$
\operatorname{sign}(e_i)=\operatorname{sign}(r_i-1), \, \forall i\in S.
$
Hence, the sign of $e_i$ retains the interpretation of over- and under-assessment. 

This identity gives a direct log-scale counterpart to the price-related dependence above, using 
$
\Cov(e,y)=\Cov(\log r,\log P).
$
Accordingly, penalizing $\Cov(e,y)$ suppresses systematic dependence between log-ratios and log-sale prices. This is well suited for the learning problem we solve for two reasons. First, it operates on the same log-scale on which the predictive model is estimated. Second, it converts multiplicative deviations in price space, represented by the ratios $r_i$, into additive residuals through the identity $e_i=\log r_i$, while also reducing the leverage of the upper tail of the price distribution.


% \medskip
% \noindent\textbf{A population intuition under squared-error prediction.}
% The dependence view above is also consistent with a basic population property of squared-error prediction in log-price space. Let $P$ denote market value, let $Y:=\log P$, and let $X$ denote the observed feature vector. Under squared-error loss, the Bayes predictor is the conditional mean
% \[
% f^\star(X)=\E[Y\mid X].
% \]
% Define the associated population log-residual by
% \[
% e^\star:=f^\star(X)-Y.
% \]
% Then
% \begin{equation}
% \label{eq:covariance_population_identity}
% \Cov(e^\star,Y)
% =
% -\E\!\left[\Var(Y\mid X)\right]
% \le 0.
% \end{equation}
% Indeed,
% \[
% \Cov(e^\star,Y)
% =
% \Cov(\E[Y\mid X]-Y,\;Y)
% =
% \Var(\E[Y\mid X])-\Var(Y),
% \]
% and the law of total variance gives
% \[
% \Var(Y)=\Var(\E[Y\mid X])+\E[\Var(Y\mid X)].
% \]
% Therefore \eqref{eq:covariance_population_identity} follows.

% This identity provides a population intuition for regressive pressure under accuracy-oriented prediction. When relevant price variation remains unobserved or imperfectly measured, conditional-mean prediction shrinks outcomes toward the center of the conditional distribution. On the log scale, this induces a non-positive association between log-residuals and log-prices, and therefore a tendency for multiplicative over-assessment of lower-valued realizations and under-assessment of higher-valued realizations. We use this result only as motivation for the dependence-based correction developed below: it does not imply that every finite-sample model trained for predictive accuracy must be regressive under every assessor metric.


\medskip
\noindent\textbf{Regressivity under squared-error prediction.} A covariance penalization in log-space further allows us to correct regressive patterns that naturally arise in learning problems, like (\ref{eq:baseline_erm}), when minimizing squared error. These patterns are characterized in Proposition \ref{prop:population_covariance_identity}.

\begin{proposition}[Population covariance sign under squared-error prediction]
\label{prop:population_covariance_identity}
Let $Y:=\log P$, let $X$ denote the observed feature vector, and let
$
f^\star(X)=\E[Y\mid X]
$
be the Bayes predictor under squared-error loss\todo{add reference?}. Define the associated population residual as
$
e^\star:=f^\star(X)-Y.
$
Then
\vspace{-2mm}
\begin{equation}
\label{eq:covariance_population_identity}
\Cov(e^\star,Y)
=
-\E\!\left[\Var(Y\mid X)\right]
\le 0.
% \vspace{-2mm}
\end{equation}
\end{proposition}

\begin{proof}
Using $f^\star(X)=\E[Y\mid X]$, we get
$
\Cov(e^\star,Y)
=
\Cov(\E[Y\mid X]-Y,\;Y)
=
\Cov(\E[Y\mid X],Y)-\Var(Y).
$
By the tower property\todo{add reference?}, we obtain
\[
\Cov(\E[Y\mid X],Y)
=
\E\!\big[\E[Y\mid X]\,Y\big]-\E[\E[Y\mid X]]\,\E[Y]
=
\E\!\big[\E[Y\mid X]^2\big]-\E[Y]^2
=
\Var(\E[Y\mid X]).
\]
Therefore,
$
\Cov(e^\star,Y)=\Var(\E[Y\mid X])-\Var(Y).
$
Finally, applying the law of total variance,
\[
\Var(Y)\geq 0  \iff \Var(\E[Y\mid X])+\E[\Var(Y\mid X)] \geq 0,
\]
which implies
$
\Cov(e^\star,Y)
=
-\E\!\left[\Var(Y\mid X)\right]
\le 0.
$
\end{proof}
% Let $Y:=\log P$, let $X$ denote the observed feature vector, and let
% $
% f^\star(X)=\E[Y\mid X]
% $
% be the \textit{Bayes predictor} under squared-error loss. Define the associated population residual vector $
% e^\star:=f^\star(X)-Y.
% $
% Then
% \begin{equation}
% \label{eq:covariance_population_identity}
% \Cov(e^\star,Y)
% =
% -\E\!\left[\Var(Y\mid X)\right]
% \le 0.
% \end{equation}
% Indeed,
% \[
% \Cov(e^\star,Y)
% =
% \Cov(\E[Y\mid X]-Y,\;Y)
% =
% \Var(\E[Y\mid X])-\Var(Y),
% \]
% and the law of total variance gives
% \[
% \Var(Y)=\Var(\E[Y\mid X])+\E[\Var(Y\mid X)].
% \]
% Therefore \eqref{eq:covariance_population_identity} follows.

Equation \eqref{eq:covariance_population_identity} shows that, whenever some conditional variation in prices remains unexplained (as it is expected in practice), squared-error prediction induces a non-positive dependence between log-residuals and log-prices. That is, accuracy-oriented regression can naturally create regressive patterns. Note that Proposition \ref{prop:cov_signs} helps us characterize a population tendency of squared-loss prediction and explain why penalizing $\Cov(e,y)$ is a natural way to counteract that pattern, but it does not imply that every finite-sample model trained to minimize squared errors will be regressive under assessors metrics defined in \ref{subsec:equity_metrics}. 

% Near the proportional-assessment benchmark $r=1$, we additionally have the first-order approximation
% \[
% \log r = (r-1) + O\!\left((r-1)^2\right),
% \]
% so shrinking $\Cov(e,y)$ also suppresses the leading-order value-related drift in the ratios around neutrality. For these reasons, $\Cov(e,y)$ serves as a natural operational surrogate for the same underlying regressivity pattern captured conceptually by $\Cov(r,P)$.

}
% In later sections, however, we work on the model's native log-price scale. Let
% \[
% y:=\log P,
% \qquad
% \hat y:=\log \widehat P,
% \qquad
% e:=\hat y-y=\log \widehat P-\log P=\log r.
% \]
% Then
% \[
% \Cov(e,y)=\Cov(\log r,\log P),
% \]
% which is the log-scale analogue of the price-related dependence above. This formulation is convenient for optimization because the model is trained in log-price space, and because the log transformation reduces the leverage of the upper tail of the price distribution. Moreover, near the neutrality benchmark $r=1$, we have $\log r \approx r-1$, so shrinking $\Cov(e,y)$ suppresses the leading-order value-related drift in the assessment ratios. This makes $\Cov(e,y)$ a natural operational surrogate for correcting regressivity while preserving the conceptual link with $\Cov(r,P)$.
% }


% \medskip

% \noindent \textbf{Log-residuals as a ratio-proxy.} Note that a key identity is obtained by comparing the residuals of the predictions in log-scale and the assessment ratios. Define the log-scale residual $
% e_i := \hat y_i-y_i
% $, which is the residual used for (\ref{eq:baseline_erm}).
% Then
% \begin{equation}
% \label{eq:log_ratio_identity}
% e_i
% =
% \log \hat P_i-\log P_i
% =
% \log\!\left(\frac{\hat P_i}{P_i}\right)= \log r_i, \qquad \forall i \in S.
% \end{equation}
% Thus, the log-space residual is exactly the logarithm of the assessment ratio. 

% \medskip

%  OLD VERSION
% \subsection{Regressivity as a Dependence Problem}
% \label{subsec:dependence_problem}

% The assessor literature usually describes regressivity as a \emph{price-related} pattern in the ratios: as property value increases, assessment ratios tend to fall. This behavior can be rapidly identified in Figure \ref{fig:regressivity_motivation_boxed}, from the negative trend between the ratios $\{r_i\}_{i\in S}$ and the log-prices $\{y_i\}_{i \in S}$; where the log-scale simplifies the example illustration.  

% \begin{figure}[!h]
%     \centering 
%     \includegraphics[width=0.7\linewidth]{Fairness Price/CCAO meeting/img/motivation/Baseline LightGBM_0_motivation_squared_only.pdf}
%     \caption{Illustration of regressive behavior in mass appraisal. Assessment-to-sale ratios $r$ tend to decrease with property value $P$: lower-value properties are more likely to be over-assessed ($r > 1$), while higher-value properties are more likely to be under-assessed ($r < 1$).}
%     \label{fig:regressivity_motivation_boxed}
% \end{figure}

% Since both the ratios and the price (or log-prices), are continuous quantities, a natural way to measure simple dependence between them is the covariance, defined as
% \begin{equation}
% \label{eq:sample_cov}
% {\operatorname{Cov}}(r,P)
% :=
% \frac{1}{|S|}
% \sum_{i\in S}
% \big(r_i-\bar r\big)\big(P_i-\bar P\big),
% \end{equation}
% where $\bar r$ and $\bar P$ denote the sample means of $\{r_i\}_{i\in S}$ and $\{P_i\}_{i\in S}$, respectively. In particular, the sign of the covariance signals the type of linear relationships between two vectors, as
% \begin{equation}
% \operatorname{\textbf{sign}}\left(\Cov (r, P)\right ) = \operatorname{\textbf{sign}}\left(\mathrm{Corr}(r,P) \right) =\operatorname{\textbf{sign}}\left(\mathrm{slope}(r \sim P) \right),
% \end{equation}
% where $\operatorname{\textbf{sign}}(x)\in \{-1,0,1\}$ is the signum function for some $x\in \R$, $\mathrm{Corr}(r,P):=\frac{\Cov (r, P)}{\sqrt{\mathrm{Var} (r) \mathrm{P} (r)}}$ denotes the correlation statistic, and $\mathrm{slope}(r\sim P):=\frac{\Cov (r, P)}{{ \mathrm{P} (r)}}$ denotes the slope obtained from a simple regression of $r$ as a function of $P$. That is, the correlation and the linear trend between $r$ and $P$ are negative (positive) if and only if the covariance between them is negative (positive). In particular, if $\Cov(r, P)<0$, then we have a negative linear dependence of $r$ in terms of $P$, which as seen in Figure \ref{fig:regressivity_motivation_boxed}, is a strong indicator of regressivity.  


% Moreover, the covariance follows a similar type of ordering relationship with some of the regressivity measures defined by the IAAO \cite{IAAO2025ExposureRatio}, introduced in $\ref{subsec:equity_metrics}$. See, e.g., Proposition \ref{prop:cov_prd}.

% \begin{proposition}
% \label{prop:cov_prd}
%     Given a random vector of the population prices $P \in \R_{++}$ and their respective assessment ratios $r = \frac{\widehat P}{P}$. The population covariance between $r$ and $P$, and the population $\mathrm{PRD} = \frac{\Expected{r}\Expected{P}}{\Expected{rP}}$ (\ref{eq:prd})
%     % assessment ratio $$
%      satisfy the ordering relationship
%      \vspace{-4mm}
%     \begin{equation}
% \mathrm{\textbf{sign}}\left(\Cov (r, P)\right ) =  \mathrm{\textbf{sign}}\left(\mathrm{1 - PRD} \right). \vspace{-4mm}
%     \end{equation}
% \end{proposition}
% \begin{proof}
%     Without loss of generality, we prove just one of the three possible relationships. Assume $\Cov(r,P) < 0$, then $\Expected{r}\Expected{P} - \Expected{r}\Expected{P} < 0$, and thus $\frac{\Expected{r}\Expected{P}}{\Expected{rP}} = PRD > 1$.  
% \end{proof}
% \noindent This means that covariance is not just a flexible indicator of regressivity, but that it is a direct diagnostic measure of regressivity as defined by the IAAO.

% % This negative  For simplicity, it is useful to rewrite this statement in a form that is directly compatible with regression outputs.

% \todo{Tal about the logaritmic scale and finalize with the remark that ML enforces regressivity.}

% By \eqref{eq:log_ratio_identity}, the log assessment ratio is exactly the log-scale residual. Hence, if we let $
% e_i := \log r_i
% \text{ and }
% y_i := \log p_i,
% $
% then any systematic price-related pattern in the ratios can equivalently be studied as a relationship between residuals and log prices. A convenient sample summary of this relationship on an evaluation sample $S$ is
% \begin{equation}
% \label{eq:sample_cov}
% \widehat{\operatorname{Cov}}_S(e,y)
% :=
% \frac{1}{|S|}
% \sum_{i\in S}
% \big(e_i-\bar e\big)\big(y_i-\bar y\big),
% \end{equation}
% where $\bar e$ and $\bar y$ denote the sample means of $\{e_i\}_{i\in S}$ and $\{y_i\}_{i\in S}$, respectively. The corresponding slope of log ratios on log prices is then
% % \begin{equation}
% % \label{eq:sample_slope}
% $
% \widehat{\beta}^{\,\log}_S
% :=
% \frac{\widehat{\operatorname{Cov}}_S(e,y)}
% {\widehat{\operatorname{Var}}_S(y)}.
% % \end{equation}
% $
% Because $\widehat{\operatorname{Var}}_S(y)$ is fixed once the evaluation sample is fixed, controlling the sign and magnitude of price-related linear trends is closely tied to controlling the covariance between log ratios and log prices.

% This training-oriented summary is \emph{not} intended to replace assessor metrics such as PRD or PRB. Rather, it serves a different role: it provides a smooth, model-compatible summary of price-related dependence that can be acted on during learning. In this sense, the paper moves from \emph{diagnosing} regressivity through ratio-study statistics to \emph{controlling} a closely related dependence pattern during training.

% \medskip
% \noindent\textbf{A population intuition under squared-error prediction.} Accuracy-oriented regression ---mainly, squared-error minimization---naturally tends to compress the extremes of the value distribution when important price drivers are unobserved or imperfectly measured. In mass appraisal, this is a mechanism that can translate into regressive behaviors.  Let $P$ denote the random market value, let $Y:=\log P$, and let $X$ denote the observed feature vector. Under squared-error loss, the Bayes\todo{maybe add some reference to the Bayes pred. and its importance?} predictor is the conditional mean
% $
% f^\star(X)=\mathbb{E}[Y\mid X].
% $
% If we define the associated residual by
% $
% e^\star := f^\star(X)-Y,
% $
% then by the law of total variance \todo{either add citation or the proof?}
% \begin{equation}
% \label{eq:covariance_population_identity}
% \operatorname{Cov}(e^\star,Y)
% =
% -\mathbb{E}\!\left[\operatorname{Var}(Y\mid X)\right]
% \le 0.
% \end{equation}
% Thus, whenever some irreducible conditional variation in prices remains, the conditional-mean predictor induces a non-positive association between log assessment ratios and log prices. Note that this identity should be interpreted just as an intuition, and not as a claim that every finite-sample model trained for predictive accuracy must be regressive under every assessor metric. 








\subsection{Temporal Evaluation and Validation Setup}
\label{subsec:oos_protocol}

\begin{figure}[!htp]
    \centering
    \includegraphics[width=0.5\linewidth]{Fairness Price/GE Paper/img/model-selection/extreme_case_3_constraint_O1_0.90_edit.pdf}
    \caption{Caption}
    \label{fig:placeholder}
\end{figure}

Mass appraisal is a temporally deployed prediction problem: models are trained on historical sales and then used to support later assessments. Hence, the notion of out-of-sample performance is chronological rather than purely random split-based. IAAO explicitly recommends holdout-based model testing for objective performance assessment \cite{IAAO2025Mass}. Further, rolling-origin evaluation is standard in the forecasting literature when model development must respect time order \cite{tashman2000out}. In this section we formalize the general evaluation procedure.

Let the full sold-data sample be ordered by sale date and partitioned into two chronologically separated parts:
$
\mathcal{D}^{\mathrm{dev}}
\quad\text{and}\quad
\mathcal{D}^{\mathrm{test}},
$
where $\mathcal{D}^{\mathrm{dev}}$ is the \textit{development} sample used for training and calibration, and $\mathcal{D}^{\mathrm{test}}$ is a final holdout \textit{test} period used only for one-time reporting after model selection is complete. 
% In other words, no model-selection decision is allowed to use information from $\mathcal{D}^{\mathrm{test}}$.

Within the development sample $\D^\mathrm{dev}$, we construct a sequence of rolling-origin windows 
$\{(\mathcal{T}_w,\mathcal{V}_w)\}_{w=1}^W,$
where each training set $\mathcal{T}_w$ precedes its corresponding validation set $\mathcal{V}_w$ in time (sale date). The windows may be expanding or otherwise office-specific, but they always respect the chronological requirement
$
\max_{i\in \mathcal{T}_w} t_i
<
\min_{j\in \mathcal{V}_w} t_j.
$
A \emph{candidate configuration} $c\in\mathcal{C}$ consists of a model class together with its hyperparameter tuning choices, including any regressivity-aware penalty parameters introduced later. For each $c\in\mathcal{C}$ and each window $w \in [W]$, we fit the model on $\mathcal{T}_w$, predict on $\mathcal{V}_w$, and record the corresponding out-of-sample performance.

To keep the notation flexible, let
$
\mathbf{A}_w(c)
\text{ and }
\mathbf{Q}_w(c)
$
denote the vectors of accuracy and equity summaries, respectively, computed on the validation sample $\mathcal{V}_w$. Accuracy metrics may be evaluated on the prediction scale or the original scale. Equity metrics are always computed after back-transforming predictions to the original price scale and forming the ratios in \eqref{eq:ratio_def}. This distinction is important and will be maintained throughout the paper: models are trained on log prices, but assessor-facing fairness is judged on price-scale ratios.

% In this paper, \emph{out-of-sample} means evaluation on sales whose dates are strictly later than the dates of the sales used to fit the corresponding model. Section~\ref{sec:model_selection} uses the collection
% \[
% \Big\{ \mathbf{A}_w(c),\mathbf{Q}_w(c)\Big\}_{c\in\mathcal{C},\,w=1,\dots,W}
% \]
% to define the multi-objective calibration and model-selection rule. After that rule selects a final configuration, the model is refit on the full development sample $\mathcal{D}^{\mathrm{dev}}$ and evaluated once on the untouched holdout sample $\mathcal{D}^{\mathrm{test}}$.

% The purpose of the present section is therefore twofold. First, it defines the ratio-study objects through which regressivity and assessment equity are measured. Second, it defines the temporal evaluation environment in which competing models will be compared. Section~\ref{sec:methodology} then uses this setup to introduce regressivity-aware regularization terms for standard regression learners.








\section{General Penalized Learning Framework}

We now formalize the in-processing approach introduced in the previous section. Recall the goal is to augment a standard predictive learner with an additional term that discourages systematic regressive distortions in the predictions, while preserving the same structure of the underlying regression model.

Let $\D := \{(x_i,y_i)\}_{i=1}^n$ be a training sample, where $x_i \in \X$ denotes the feature vector of property $i \in [n]$ and $y_i \in \R$ denotes its log sale price. Let $\Theta$ be a parameters space, $f_\theta:\X \to \R$ be the prediction function, for each $\theta \in \Theta$. As in (\ref{eq:baseline_erm}), the \textit{baseline} predictive model is defined as $f_{\theta_\mathrm{base}}$, where
\begin{align}
    \label{opt:base_model}
    \theta_{\mathrm{base}}
    \in
    \arg\min_{\theta \in \Theta}
    \left\{
        L_n(f_\theta) + \Omega
    \right\},
    \qquad
    L_n(f_\theta) := \frac{1}{n}\sum_{i=1}^n \ell(f_\theta(x_i),y_i),
\end{align}
where $\ell$ is a predefined regression loss function (e.g., squared-error) and $\Omega:\Theta \to \R$ is an optional complexity regularizer. Throughout this section, all quantities are empirical unless otherwise stated.

To connect this predictive problem with assessment equity, define the fitted log value
\[
\hat y_i(f) := f(x_i),
\]
and the log-scale residual
\[
e_i(f) := \hat y_i(f) - y_i.
\]
Since $y_i = \log P_i$ and $\hat y_i(f) = \log \hat P_i(f)$, where $P_i$ and $\hat P_i(f)$ denote the market value and fitted assessed value of property $i$, respectively, we have
\begin{align}
    % \label{eq:log_ratio_identity}
    e_i(f)
    = \log \hat P_i(f) - \log P_i
    = \log\!\left(\frac{\hat P_i(f)}{P_i}\right)
    = \log r_i(f),
\end{align}
where
\[
r_i(f) := \frac{\hat P_i(f)}{P_i}
\]
is the assessment-to-sale ratio induced by the predictor $f$. Thus, controlling systematic dependence between the residuals $e_i(f)$ and the log values $y_i$ is equivalent to controlling systematic dependence between log assessment ratios and log market values.

Motivated by this equivalence, let $\Psi_n(f)$ denote a sample-dependent functional designed to measure regressivity-relevant dependence in the fitted assessments. One may then formulate a constrained learning problem of the form
\begin{align}
    \label{opt:general_constrained_model}
    \min_{f \in \H}
    \left\{
        L_n(f) + \mathcal{R}(f)
    \right\}
    \quad
    \text{s.t.}
    \quad
    \Psi_n(f) \le \tau,
\end{align}
for some tolerance level $\tau \ge 0$. In this paper, we work instead with the corresponding soft-constrained penalized formulation
\begin{align}
    \label{opt:general_penalized_model}
    f_{\rho}
    \in
    \arg\min_{f \in \H}
    \left\{
        L_n(f) + \mathcal{R}(f) + \rho\,\Psi_n(f)
    \right\},
\end{align}
where $\rho \ge 0$ is a penalty parameter that controls the strength of the regressivity correction. The baseline model \eqref{opt:base_model} is recovered as the special case $\rho=0$.

Formulation \eqref{opt:general_penalized_model} isolates the general structure of the proposed approach: starting from a standard predictive learner, we augment the empirical objective with a penalty that discourages systematic value-dependent distortions in the resulting assessments. In the present setting, the relevant distortions are naturally expressed through the dependence between assessment ratios and property values, or equivalently, between log-ratios and log-prices. The next section makes this idea concrete by introducing covariance-based choices of the functional $\Psi_n$, including both direct covariance penalties and computationally convenient surrogate forms.

\paragraph{Remark.}
Once a concrete choice of $\Psi_n$ is specified, problem \eqref{opt:general_penalized_model} can be instantiated in standard learners used in mass appraisal, including regularized linear models and tree-based boosting methods. These model-specific implementations are developed after the covariance-based penalties have been introduced.

\section{Covariance-Based Penalizations}

We now instantiate the general penalized learning problem \eqref{opt:general_penalized_model} by specifying concrete choices of the functional $\Psi_n(f)$ based on empirical covariance-type measures. Throughout, we keep the notation of the previous section: for a predictor $f \in \H$,
\[
\hat y_i(f) = f(x_i),
\qquad
e_i(f) = \hat y_i(f) - y_i,
\qquad
r_i(f) = \exp(e_i(f)),
\]
so that $e_i(f)=\log r_i(f)$ by \eqref{eq:log_ratio_identity}. Since regressivity corresponds to a systematic decrease of assessment ratios with property values, a natural strategy is to penalize dependence between the vectors $e(f):=(e_1(f),\dots,e_n(f))$ and $y:=(y_1,\dots,y_n)$.

Let
\[
\bar e(f) := \frac{1}{n}\sum_{i=1}^n e_i(f),
\qquad
\bar y := \frac{1}{n}\sum_{i=1}^n y_i,
\]
and define the empirical covariance
\begin{align}
    \label{eq:empirical_covariance}
    \widehat{\Cov}_n(e(f),y)
    :=
    \frac{1}{n}\sum_{i=1}^n
    \bigl(e_i(f)-\bar e(f)\bigr)\bigl(y_i-\bar y\bigr).
\end{align}
This quantity is directly connected to the slope of the least-squares regression of $e_i(f)$ on $y_i$, since
\[
\widehat{\beta}_{e\sim y}(f)
=
\frac{\widehat{\Cov}_n(e(f),y)}{\widehat{\Var}_n(y)},
\]
and $\widehat{\Var}_n(y)$ does not depend on the model choice. Accordingly, controlling the empirical covariance is a natural way to control first-order value-dependent trends in the fitted assessment ratios.

\subsection{Direct covariance penalties}

A first class of penalties is obtained by using the empirical covariance in \eqref{eq:empirical_covariance} directly.

\paragraph{Squared covariance penalty.}
The most symmetric choice is to penalize the squared covariance:
\begin{align}
    \label{eq:cov2_penalty}
    \Psi_n^{\mathrm{cov}^2}(f)
    :=
    \frac{1}{2}\,
    \widehat{\Cov}_n(e(f),y)^2.
\end{align}
Substituting \eqref{eq:cov2_penalty} into \eqref{opt:general_penalized_model} yields
\begin{align}
    \label{opt:cov2_penalized_model}
    f_{\rho}^{\mathrm{cov}^2}
    \in
    \arg\min_{f\in\H}
    \left\{
        L_n(f)+\mathcal{R}(f)
        +
        \frac{\rho}{2}\,\widehat{\Cov}_n(e(f),y)^2
    \right\}.
\end{align}
This choice penalizes any linear dependence between $e(f)$ and $y$, regardless of sign. In particular, it suppresses both regressive and progressive linear trends, and is therefore appropriate when the goal is to drive the first-order dependence toward zero.

\paragraph{Signed covariance penalty.}
In some applications, one may wish to focus specifically on regressive patterns under the sign convention adopted here, namely $e_i(f)=\log r_i(f)$. Since regressivity corresponds to negative dependence between $e(f)$ and $y$, a directional alternative is
\begin{align}
    \label{eq:negcov_penalty}
    \Psi_n^{\mathrm{neg\mbox{-}cov}}(f)
    :=
    -\,\widehat{\Cov}_n(e(f),y).
\end{align}
The resulting optimization problem is
\begin{align}
    \label{opt:negcov_penalized_model}
    f_{\rho}^{\mathrm{neg\mbox{-}cov}}
    \in
    \arg\min_{f\in\H}
    \left\{
        L_n(f)+\mathcal{R}(f)
        -
        \rho\,\widehat{\Cov}_n(e(f),y)
    \right\}.
\end{align}
Unlike \eqref{opt:cov2_penalized_model}, this formulation is directional rather than symmetric: it directly pushes the fitted model away from negative covariance and toward less regressive dependence structures. We therefore view \eqref{opt:negcov_penalized_model} as a sign-sensitive alternative when only regressive trends, rather than all linear trends, are the primary concern.

\subsection{A separable surrogate}

While the direct covariance penalties above are conceptually natural, they have two important limitations. First, covariance is a first-order summary and may fail to reflect more complex dependence patterns when positive and negative local trends cancel each other out. Second, the term $\widehat{\Cov}_n(e(f),y)^2$ is non-separable across observations, which can be inconvenient in optimization methods that rely on observation-wise gradients or second-order approximations.

To motivate a more computationally convenient alternative, define
\[
a_i(f):=\bigl(e_i(f)-\bar e(f)\bigr)\bigl(y_i-\bar y\bigr).
\]
Then,
\[
\widehat{\Cov}_n(e(f),y)
=
\frac{1}{n}\sum_{i=1}^n a_i(f),
\]
and Jensen's inequality implies
\begin{align}
    \label{eq:covariance_upper_bound}
    \widehat{\Cov}_n(e(f),y)^2
    =
    \left(
        \frac{1}{n}\sum_{i=1}^n a_i(f)
    \right)^2
    \le
    \frac{1}{n}\sum_{i=1}^n a_i(f)^2
    =
    \frac{1}{n}\sum_{i=1}^n
    \bigl(e_i(f)-\bar e(f)\bigr)^2
    \bigl(y_i-\bar y\bigr)^2.
\end{align}
The right-hand side of \eqref{eq:covariance_upper_bound} is an upper bound on the squared covariance, but it still depends on the sample mean $\bar e(f)$ and therefore remains coupled across observations.

In many well-calibrated regression settings, however, the average residual is expected to be small, and perfect equity corresponds to $e_i(f)=0$ for every $i$, or equivalently $r_i(f)=1$. Motivated by this observation, we consider the simplified surrogate
\begin{align}
    \label{eq:surrogate_penalty}
    \Psi_n^{\mathrm{surr}}(f)
    :=
    \frac{1}{n}\sum_{i=1}^n
    e_i(f)^2\bigl(y_i-\bar y\bigr)^2.
\end{align}
Substituting \eqref{eq:surrogate_penalty} into \eqref{opt:general_penalized_model} gives the surrogate-penalized problem
\begin{align}
    \label{opt:surrogate_penalized_model}
    f_{\rho}^{\mathrm{surr}}
    \in
    \arg\min_{f\in\H}
    \left\{
        L_n(f)+\mathcal{R}(f)
        +
        \rho\,\Psi_n^{\mathrm{surr}}(f)
    \right\}.
\end{align}

The surrogate \eqref{eq:surrogate_penalty} has two useful features. First, it remains anchored to the same underlying idea as the covariance penalties: residual distortions are weighted more heavily when they occur far from the center of the value distribution. Second, unlike \eqref{eq:cov2_penalty}, it is separable across observations once the weights $(y_i-\bar y)^2$ are fixed. This makes \eqref{opt:surrogate_penalized_model} particularly attractive for model classes and optimization procedures in which observation-wise objective contributions are computationally convenient.

\subsection{Discussion}

The three formulations above target related but distinct notions of value-dependent distortion. The squared covariance penalty \eqref{eq:cov2_penalty} is the most direct way to suppress linear dependence between log-ratios and log-values. The signed covariance penalty \eqref{eq:negcov_penalty} focuses instead on the direction of dependence associated with regressivity under the present sign convention. The surrogate penalty \eqref{eq:surrogate_penalty} sacrifices exact covariance control in exchange for a data-point-wise separable objective that is better suited to some learning algorithms.

In the remainder of the paper, these penalties are instantiated within standard regression models used in practice by assessors. In particular, the next section shows how the covariance-based objectives above can be implemented in model classes such as regularized linear regression and gradient boosting, thereby preserving the general learning framework while yielding practically usable regressivity-aware predictors.

\section{Model-Specific Implementations}
\label{sec:model_specific_implementations}

We now show how the covariance-based penalizations introduced in the previous section can be instantiated in specific model classes. We focus on two examples that are especially relevant in practice: linear models and second-order tree boosting. The former provides a transparent analytical benchmark, while the latter covers widely used high-performing nonlinear learners such as XGBoost- and LightGBM-type methods.

\subsection{Linear models}
\label{subsec:linear_models}

Consider the linear hypothesis class
\[
\H_{\mathrm{lin}} := \{f_\theta : x \mapsto x^\top \theta \;:\; \theta \in \Theta \subseteq \R^p\},
\]
and let $X \in \R^{n\times p}$ denote the design matrix. For a coefficient vector $\theta \in \Theta$, define
\[
f_\theta(x_i) = x_i^\top \theta,
\qquad
e := X\theta - y.
\]
As in the previous section, let
\[
H := I - \frac{1}{n}\mathbf{1}\mathbf{1}^\top,
\qquad
c := Hy = y - \bar y\,\mathbf{1},
\qquad
D := \diag(c_1^2,\dots,c_n^2),
\]
so that
\[
\widehat{\Cov}_n(e,y)
=
\frac{1}{n}e^\top c.
\]

Throughout this subsection, we take the baseline loss to be the squared loss
\[
L_n(f_\theta) = \frac{1}{2n}\|X\theta-y\|_2^2.
\]
We then derive the corresponding penalized estimators for the three covariance-based choices introduced in Section~\ref{sec:covariance_penalizations}.

\paragraph{Signed covariance penalty.}
For the directional penalty
\[
\Psi_n^{\mathrm{neg\mbox{-}cov}}(f_\theta)
=
-\widehat{\Cov}_n(e,y)
=
-\frac{1}{n}e^\top c,
\]
the penalized objective is
\begin{align}
    \label{opt:linear_negcov}
    \theta_{\rho}^{\mathrm{neg\mbox{-}cov}}
    \in
    \arg\min_{\theta\in\Theta}
    \left\{
        \frac{1}{2n}\|X\theta-y\|_2^2
        -
        \frac{\rho}{n}e^\top c
    \right\}.
\end{align}
Completing the square gives
\[
\frac{1}{2n}\|X\theta-y\|_2^2
-
\frac{\rho}{n}e^\top c
=
\frac{1}{2n}\|X\theta-(y+\rho c)\|_2^2
-
\frac{\rho^2}{2n}\|c\|_2^2,
\]
so \eqref{opt:linear_negcov} is equivalent, up to an additive constant, to ordinary least squares with shifted response $y+\rho c$. Consequently, whenever $X^\top X$ is invertible,
\begin{align}
    \label{eq:linear_negcov_closed_form}
    \theta_{\rho}^{\mathrm{neg\mbox{-}cov}}
    =
    (X^\top X)^{-1}X^\top (y+\rho c).
\end{align}
Thus, under squared loss, the signed covariance penalty is particularly simple to implement in linear models.

\paragraph{Squared covariance penalty.}
For the symmetric penalty
\[
\Psi_n^{\mathrm{cov}^2}(f_\theta)
=
\frac{1}{2}\,\widehat{\Cov}_n(e,y)^2
=
\frac{1}{2n^2}\bigl(e^\top c\bigr)^2,
\]
the penalized estimator solves
\begin{align}
    \label{opt:linear_cov2}
    \theta_{\rho}^{\mathrm{cov}^2}
    \in
    \arg\min_{\theta\in\Theta}
    \left\{
        \frac{1}{2n}\|X\theta-y\|_2^2
        +
        \frac{\rho}{2n^2}\bigl((X\theta-y)^\top c\bigr)^2
    \right\}.
\end{align}
Define
\[
a := X^\top c,
\qquad
b := y^\top c = c^\top c.
\]
Then the first-order optimality condition for \eqref{opt:linear_cov2} is
\[
\left(
\frac{1}{n}X^\top X + \frac{\rho}{n^2}aa^\top
\right)\theta
=
\frac{1}{n}X^\top y + \frac{\rho}{n^2}ba.
\]
Hence, whenever the displayed inverse exists,
\begin{align}
    \label{eq:linear_cov2_closed_form}
    \theta_{\rho}^{\mathrm{cov}^2}
    =
    \left(
        \frac{1}{n}X^\top X
        +
        \frac{\rho}{n^2}aa^\top
    \right)^{-1}
    \left(
        \frac{1}{n}X^\top y
        +
        \frac{\rho}{n^2}ba
    \right).
\end{align}
This shows that the squared covariance penalty induces a rank-one modification of the usual normal equations.

\paragraph{Separable surrogate penalty.}
For the separable surrogate
\[
\Psi_n^{\mathrm{surr}}(f_\theta)
=
\frac{1}{n}\sum_{i=1}^n e_i^2 c_i^2
=
\frac{1}{n}e^\top D\,e,
\]
the penalized problem is
\begin{align}
    \label{opt:linear_surr}
    \theta_{\rho}^{\mathrm{surr}}
    \in
    \arg\min_{\theta\in\Theta}
    \left\{
        \frac{1}{2n}\|X\theta-y\|_2^2
        +
        \frac{\rho}{n}e^\top D\,e
    \right\}.
\end{align}
The corresponding optimality condition is
\[
\left(
\frac{1}{n}X^\top X
+
\frac{2\rho}{n}X^\top D X
\right)\theta
=
\frac{1}{n}X^\top y
+
\frac{2\rho}{n}X^\top D y,
\]
and therefore, whenever the displayed inverse exists,
\begin{align}
    \label{eq:linear_surr_closed_form}
    \theta_{\rho}^{\mathrm{surr}}
    =
    \left(
        \frac{1}{n}X^\top X
        +
        \frac{2\rho}{n}X^\top D X
    \right)^{-1}
    \left(
        \frac{1}{n}X^\top y
        +
        \frac{2\rho}{n}X^\top D y
    \right).
\end{align}
The factor \(2\rho\) arises directly from the definition of \(\Psi_n^{\mathrm{surr}}\) in Section~\ref{sec:covariance_penalizations}.

\paragraph{Extension to regularized linear models.}
The same constructions extend immediately to penalized linear estimators of the form
\[
\min_{\theta\in\Theta}
\left\{
    L_n(f_\theta) + \mathcal{R} + \rho\,\Psi_n(f_\theta)
\right\},
\]
where \(\mathcal{R}\) is a complexity regularizer on the coefficients. In the ridge case, the resulting systems remain quadratic and admit closed-form solutions after adding the usual ridge term to the coefficient matrix. In contrast, \(\ell_1\)-regularized variants such as the LASSO remain convex but no longer admit closed-form expressions, and must instead be solved numerically.

\subsection{Second-order Gradient boosting}
\label{subsec:tree_boosting}

We next consider second-order additive tree boosting under squared loss, as used in gradient-boosted tree systems such as XGBoost and LightGBM. Let \(f^{(m-1)}\) denote the current prediction function at boosting iteration \(m\), and write
\[
f_i := f^{(m-1)}(x_i),
\qquad
e_i := f_i - y_i,
\qquad
c_i := y_i - \bar y,
\qquad
z := \widehat{\Cov}_n(e,y) = \frac{1}{n}\sum_{i=1}^n e_i c_i.
\]
At step \(m\), one adds a regression tree \(t\) to the current predictor, so that
\[
f^{(m)}(x) = f^{(m-1)}(x) + \nu\,t(x),
\]
where \(\nu\in(0,1]\) is the learning rate.

For standard second-order boosting, the objective is approximated locally by
\begin{align}
    \label{eq:boosting_local_objective}
    \widetilde{\mathcal J}(t)
    =
    \sum_{i=1}^n
    \left(
        g_i\,t_i + \frac{1}{2}h_i\,t_i^2
    \right)
    +
    \Omega(t),
\end{align}
where \(t_i := t(x_i)\), \(\Omega(t)\) is the tree-complexity penalty, and
\[
g_i := \frac{\partial}{\partial f_i}
\left(
L_n(f) + \rho\,\Psi_n(f)
\right),
\qquad
h_i := \frac{\partial^2}{\partial f_i^2}
\left(
L_n(f) + \rho\,\Psi_n(f)
\right)
\]
are evaluated at the current predictions. Under squared loss,
\[
\frac{\partial L_n(f)}{\partial f_i} = \frac{e_i}{n},
\qquad
\frac{\partial^2 L_n(f)}{\partial f_i^2} = \frac{1}{n}.
\]
The covariance-based penalties modify these quantities as follows.

\paragraph{Signed covariance penalty.}
For \(\Psi_n^{\mathrm{neg\mbox{-}cov}}(f) = -n^{-1}\sum_{i=1}^n e_i c_i\), we obtain
\begin{align}
    \label{eq:boosting_negcov_grad_hess}
    g_i^{\mathrm{neg\mbox{-}cov}}
    =
    \frac{e_i}{n}
    -
    \frac{\rho}{n}c_i,
    \qquad
    h_i^{\mathrm{neg\mbox{-}cov}}
    =
    \frac{1}{n}.
\end{align}
Thus, the signed covariance penalty yields an exact observation-wise objective and can be implemented directly through standard custom-gradient interfaces.

\paragraph{Squared covariance penalty.}
For \(\Psi_n^{\mathrm{cov}^2}(f)=\tfrac12 z^2\), the gradient is
\begin{align}
    \label{eq:boosting_cov2_gradient}
    g_i^{\mathrm{cov}^2}
    =
    \frac{e_i}{n}
    +
    \frac{\rho}{n}z\,c_i.
\end{align}
The exact Hessian of the penalized objective with respect to the prediction vector \(f=(f_1,\dots,f_n)\) is
\begin{align}
    \label{eq:boosting_cov2_exact_hessian}
    \nabla_f^2
    \bigl(L_n(f)+\rho\,\Psi_n^{\mathrm{cov}^2}(f)\bigr)
    =
    \frac{1}{n}I
    +
    \frac{\rho}{n^2}cc^\top.
\end{align}
Unlike the separable boosting objectives used in standard libraries, \eqref{eq:boosting_cov2_exact_hessian} is dense: each observation is coupled to every other observation through the rank-one term \(cc^\top\). Consequently, there are two natural implementation routes.

The first is a practical approximation that keeps only the diagonal terms,
\begin{align}
    \label{eq:boosting_cov2_diag_hessian}
    h_i^{\mathrm{cov}^2,\mathrm{diag}}
    =
    \frac{1}{n}
    +
    \frac{\rho}{n^2}c_i^2.
\end{align}
This produces a custom objective compatible with standard second-order boosting APIs, at the price of ignoring the cross-observation curvature induced by the covariance penalty.

The second is an exact leaf-wise implementation for a fixed tree structure. Suppose a candidate tree \(t\) has leaves \(I_1,\dots,I_J\) and leaf weights \(w=(w_1,\dots,w_J)\), so that \(t_i=w_j\) for \(i\in I_j\). Let
\[
n_j := |I_j|,
\qquad
b_j := \sum_{i\in I_j} e_i,
\qquad
s_j := \sum_{i\in I_j} c_i,
\]
and let \(s=(s_1,\dots,s_J)^\top\). If \(\Omega(t)=\gamma J + \frac{\lambda}{2}\|w\|_2^2\), then the exact local objective for the squared covariance penalty can be written as
\begin{align}
    \label{eq:boosting_cov2_exact_leaf_objective}
    \widetilde{\mathcal J}_{\mathrm{cov}^2}(w)
    =
    g_T^\top w
    +
    \frac{1}{2}w^\top A_T w
    +
    \gamma J
    + \text{const},
\end{align}
where
\[
g_{T,j}
=
\frac{b_j}{n}
+
\frac{\rho}{n}z\,s_j,
\qquad
A_T
=
\diag\!\left(\frac{n_1}{n}+\lambda,\dots,\frac{n_J}{n}+\lambda\right)
+
\frac{\rho}{n^2}ss^\top.
\]
Hence the optimal leaf weights satisfy
\begin{align}
    \label{eq:boosting_cov2_exact_leaf_system}
    A_T w^\star = -g_T,
    \qquad
    w^\star = -A_T^{-1}g_T,
\end{align}
whenever \(A_T\) is invertible. The associated gain for the fixed structure \(T\) is
\begin{align}
    \label{eq:boosting_cov2_exact_gain}
    \operatorname{Gain}(T)
    =
    \frac{1}{2}g_T^\top A_T^{-1}g_T - \gamma J.
\end{align}
This exact route preserves the full curvature of the covariance penalty, but requires a custom tree-construction procedure rather than the usual observation-wise Hessian interface.

\paragraph{Separable surrogate penalty.}
For the surrogate
\[
\Psi_n^{\mathrm{surr}}(f)
=
\frac{1}{n}\sum_{i=1}^n e_i^2 c_i^2,
\]
the gradient and Hessian remain observation-wise separable:
\begin{align}
    \label{eq:boosting_surr_grad_hess}
    g_i^{\mathrm{surr}}
    &=
    \frac{e_i}{n}
    +
    \frac{2\rho}{n}e_i c_i^2,
    \\
    h_i^{\mathrm{surr}}
    &=
    \frac{1}{n}
    +
    \frac{2\rho}{n}c_i^2.
\end{align}
Accordingly, the surrogate penalty is particularly well suited to off-the-shelf second-order boosting implementations: once the weights \(c_i^2\) are precomputed, the penalized objective can be supplied through standard custom gradient and Hessian routines.

\paragraph{Implementation summary.}
The three penalties above lead to distinct computational regimes in boosted trees. The signed covariance penalty is fully separable and is therefore the simplest directional correction to implement. The surrogate penalty is also separable, while placing greater emphasis on distortions occurring away from the center of the value distribution. The squared covariance penalty is the most direct implementation of the covariance-control idea, but introduces dense curvature; in practice, it can be handled either through a diagonal Hessian approximation in standard boosting libraries or through an exact custom leaf solver of the form \eqref{eq:boosting_cov2_exact_leaf_system}.

The formulas in this section show that the covariance-based framework is compatible with both transparent linear estimators and high-performing nonlinear learners. This model-specific flexibility is important for the practical goal of the paper: regressivity correction can be layered onto standard predictive tools already used in assessment workflows, rather than requiring a separate valuation architecture.

\paragraph{Algorithmic details.}
For completeness, Appendix~\ref{app:boosting_algorithms} provides pseudocode for an exact Newton-style tree-boosting implementation of the squared covariance penalty, together with an optional first-order variant. In the main text, we focus on the gradient, Hessian, and leaf-system formulas that define the method and clarify its relation to standard boosted-tree implementations.

\section{Constrained Model Selection Across Validation Folds}
\label{subsec:model_selection}



The covariance-regularized learners developed in the previous sections enlarge the set of candidate specifications available at calibration time, but they do not alter the basic structure of the assessor's workflow. After fitting the implemented model families---namely regularized linear models and gradient boosting models, with or without covariance-based penalization---the practical task remains to compare a finite collection of candidate specifications on out-of-sample validation data and to select one final specification for deployment. The difference is that, once regressivity-aware penalties are introduced, candidate models may trade off predictive performance, ratio-study diagnostics, regressivity measures, and dispersion measures in different ways. Accordingly, model calibration is naturally formulated as a constrained multi-criteria selection problem.

\paragraph{Single-model calibration.}
Let $\mathcal V=\{1,\dots,V\}$ index the validation folds (or rolling-origin validation windows), and let $V_v$ denote the validation sample associated with fold $v\in\mathcal V$. Let $\mathcal C$ denote a finite set of candidate model configurations. Each configuration $c\in\mathcal C$ specifies
(i) a model family from the implemented classes,
(ii) a penalty choice
\[
\Psi_n \in \{0,\Psi_n^{\mathrm{cov}^2},\Psi_n^{\mathrm{neg\mbox{-}cov}},\Psi_n^{\mathrm{surr}}\},
\]
where $\Psi_n=0$ denotes the unpenalized baseline,
and
(iii) all associated tuning parameters, including the penalty level $\rho\ge 0$ when applicable.
For each fold $v\in\mathcal V$, fitting configuration $c$ on the corresponding training portion yields a predictor $f_{c,v}$.

For each observation $i\in V_v$, let
\[
\hat y_i^{\,v}(c):=f_{c,v}(x_i^{\,v}),
\qquad
\hat P_i^{\,v}(c):=\exp\!\bigl(\hat y_i^{\,v}(c)\bigr),
\qquad
r_i^{\,v}(c):=\frac{\hat P_i^{\,v}(c)}{P_i^{\,v}},
\]
where $P_i^{\,v}$ is the observed sale price. From the validation predictions on fold $v$, we compute a predictive performance score $\mathrm{Perf}_v(c)$, written so that larger values are preferred, together with a collection of fold-level assessment diagnostics $Q_{m,v}(c)$ indexed by $m$. If one starts from an error metric $E_v(c)$ that is minimized, such as MSE, RMSE, or MAE, one may simply set $\mathrm{Perf}_v(c):=-E_v(c)$.

To aggregate results across folds, let $\alpha_v\ge 0$ satisfy $\sum_{v\in\mathcal V}\alpha_v=1$, and define
\[
\bar{\mathrm{Perf}}(c):=\sum_{v\in\mathcal V}\alpha_v\,\mathrm{Perf}_v(c),
\qquad
\bar Q_m(c):=\sum_{v\in\mathcal V}\alpha_v\,Q_{m,v}(c).
\]
The default choice is uniform weighting, $\alpha_v=1/V$, although the same formulation also allows non-uniform weights when one wishes to emphasize, for example, more recent validation windows.

We then select a single deployed model by solving
\begin{align}
    \label{opt:model_selection_single}
    c^\star
    \in
    \arg\max_{c\in\mathcal C}
    \quad
    & \bar{\mathrm{Perf}}(c)
    \\
    \text{s.t.}\quad
    & \underline{\kappa}_m \le \bar Q_m(c) \le \overline{\kappa}_m,
    \qquad m\in\mathcal M_{\mathrm{int}},
    \nonumber\\
    & \bar Q_m(c) \le \overline{\kappa}_m,
    \qquad m\in\mathcal M_{\mathrm{up}}.
    \nonumber
\end{align}
Here, $\mathcal M_{\mathrm{int}}$ indexes diagnostics constrained to lie in admissible intervals, while $\mathcal M_{\mathrm{up}}$ indexes diagnostics handled through upper bounds only. Typical elements of $\mathcal M_{\mathrm{int}}$ include ratio statistics such as the mean ratio, weighted mean ratio, and median ratio, as well as regressivity diagnostics such as VEI, PRD, and PRB. Typical elements of $\mathcal M_{\mathrm{up}}$ include dispersion measures such as COD and COV. Additional assessor-facing diagnostics can be incorporated by adding them to either set according to the desired feasibility requirement.

Formulation \eqref{opt:model_selection_single} deliberately avoids collapsing all validation criteria into a single weighted score. Predictive performance remains the objective, while ratio-study, regressivity, and dispersion diagnostics enter as explicit feasibility requirements. This mirrors the logic of assessor calibration more closely than a purely accuracy-driven ranking rule: among all candidate specifications that satisfy the chosen assessment-quality requirements, select the one with the best aggregated out-of-sample predictive performance.

\paragraph{How the constrained calibration differs from standard accuracy-only selection.}
The single-model procedure can be read as the following practical recipe.
\begin{enumerate}
    \item Specify a finite candidate set $\mathcal C$ of model configurations, including baseline and covariance-penalized variants.
    \item For each $c\in\mathcal C$ and each validation fold $v\in\mathcal V$, fit the model on the corresponding training portion and compute validation predictions.
    \item On each fold, compute the chosen predictive score $\mathrm{Perf}_v(c)$ together with the selected assessment diagnostics $Q_{m,v}(c)$.
    \item Aggregate these fold-level quantities into $\bar{\mathrm{Perf}}(c)$ and $\bar Q_m(c)$.
    \item Discard all configurations that violate any active assessment-quality constraint.
    \item Rank the remaining feasible configurations by aggregated predictive performance and select the best one.
    \item Refit the selected configuration $c^\star$ on the full pre-assessment development sample and deploy that final predictor.
\end{enumerate}
The only substantive departure from the usual accuracy-only workflow lies in Steps 5--6: instead of ranking all candidates solely by predictive error, the constrained version first screens for feasibility with respect to the chosen assessor-facing diagnostics and then ranks only the feasible candidates by predictive performance.

The formulation above uses aggregated fold-level diagnostics as the default choice. When temporal stability is especially important, one may replace some aggregated constraints with foldwise requirements of the form
\[
Q_{m,v}(c)\in[\underline{\kappa}_m,\overline{\kappa}_m]
\qquad
\forall v\in\mathcal V,
\]
or, more generally, use a more conservative fold aggregator. We keep the weighted-average formulation as the main presentation because it yields a transparent calibration problem over a finite candidate set and aligns naturally with standard validation practice.

\subsubsection{Convex stacking as an extension}
\label{subsubsec:convex_stacking}

A natural extension of \eqref{opt:model_selection_single} is to allow convex combinations of candidate predictions rather than selecting only one configuration. This does not change the objective--constraint logic of the calibration step; it changes only the feasible set. In the single-model case, the decision variable is a discrete configuration \(c\in\mathcal C\), and calibration is therefore a constrained sorting problem over finitely many candidates. In the stacking case, the decision variable becomes a nonnegative weight vector on the simplex, and calibration becomes a constrained optimization problem over those weights.

Let $\mathcal C_{\mathrm{stk}}=\{c_1,\dots,c_M\}\subseteq\mathcal C$ denote the set of candidate configurations admitted into the stacking pool. For each fold $v\in\mathcal V$ and observation $i\in V_v$, let
\[
\hat P_{i,j}^{\,v}:=\hat P_i^{\,v}(c_j)
\]
denote the out-of-fold predicted assessed value on the original price scale produced by configuration $c_j$. For a weight vector
\[
w=(w_1,\dots,w_M)\in\Delta_M
:=
\left\{
w\in\mathbb R_+^M : \sum_{j=1}^M w_j = 1
\right\},
\]
define the stacked prediction on fold $v$ by
\begin{align}
    \label{eq:stacked_prediction}
    \hat P_i^{\,v}(w)
    :=
    \sum_{j=1}^M w_j\,\hat P_{i,j}^{\,v},
    \qquad
    r_i^{\,v}(w)
    :=
    \frac{\hat P_i^{\,v}(w)}{P_i^{\,v}}.
\end{align}
We emphasize that stacking is defined on the \emph{price scale}, not on the log-prediction scale. This choice is natural because the assessment diagnostics entering the constraints are themselves ratio-study quantities defined in terms of assessed values and sale prices.

Given \eqref{eq:stacked_prediction}, one can compute fold-level performance and assessment diagnostics exactly as in the single-model case, now viewed as functions of the weights:
\[
\mathrm{Perf}_v(w),
\qquad
Q_{m,v}(w).
\]
The stacked calibration problem is then
\begin{align}
    \label{opt:model_selection_stacking}
    w^\star
    \in
    \arg\max_{w\in\Delta_M}
    \quad
    & \sum_{v\in\mathcal V}\alpha_v\,\mathrm{Perf}_v(w)
    \\
    \text{s.t.}\quad
    & \underline{\kappa}_m
      \le
      \sum_{v\in\mathcal V}\alpha_v\,Q_{m,v}(w)
      \le
      \overline{\kappa}_m,
      \qquad
      m\in\mathcal M_{\mathrm{int}},
      \nonumber\\
    & \sum_{v\in\mathcal V}\alpha_v\,Q_{m,v}(w)
      \le
      \overline{\kappa}_m,
      \qquad
      m\in\mathcal M_{\mathrm{up}}.
      \nonumber
\end{align}
Thus, the stacking extension preserves the same fold structure, the same diagnostics, and the same feasibility logic as \eqref{opt:model_selection_single}; it simply replaces discrete candidate comparison by continuous optimization over convex combinations of their out-of-fold predictions.

Once an optimal weight vector $w^\star$ has been identified, the final deployed stacked predictor is obtained by refitting each constituent configuration $c_j\in\mathcal C_{\mathrm{stk}}$ on the full pre-assessment development sample and combining their resulting price-scale predictions using the weights $w^\star$.

\subsubsection{Illustrative convex stacking example}
\label{subsubsec:convex_example}

A useful special case is obtained by taking average validation MSE as the objective and imposing worst-fold box constraints on selected ratio-study diagnostics. Let \(n_v:=|V_v|\). For each fold \(v\), define the stacked MSE
\[
\mathrm{MSE}_v(w)
:=
\frac{1}{n_v}\sum_{i\in V_v}
\bigl(\hat P_i^{\,v}(w)-P_i^{\,v}\bigr)^2.
\]
Consider the problem
\begin{align}
    \label{opt:stacking_convex_example}
    \min_{w\in\Delta_M}
    \quad
    & \frac{1}{V}\sum_{v\in\mathcal V}\mathrm{MSE}_v(w)
    \\
    \text{s.t.}\quad
    & \underline{\tau}_{\mathrm{mean}}
      \le
      \frac{1}{n_v}\sum_{i\in V_v}\frac{\hat P_i^{\,v}(w)}{P_i^{\,v}}
      \le
      \overline{\tau}_{\mathrm{mean}},
      \qquad \forall v\in\mathcal V,
      \nonumber\\
    & \underline{\tau}_{\mathrm{wmean}}
      \le
      \frac{\sum_{i\in V_v}\hat P_i^{\,v}(w)}{\sum_{i\in V_v}P_i^{\,v}}
      \le
      \overline{\tau}_{\mathrm{wmean}},
      \qquad \forall v\in\mathcal V,
      \nonumber\\
    & \underline{\tau}_{\mathrm{PRD}}
      \le
      \frac{
        \frac{1}{n_v}\sum_{i\in V_v}\hat P_i^{\,v}(w)/P_i^{\,v}
      }{
        \sum_{i\in V_v}\hat P_i^{\,v}(w)\big/\sum_{i\in V_v}P_i^{\,v}
      }
      \le
      \overline{\tau}_{\mathrm{PRD}},
      \qquad \forall v\in\mathcal V.
      \nonumber
\end{align}
The objective in \eqref{opt:stacking_convex_example} is convex quadratic in \(w\), because each \(\hat P_i^{\,v}(w)\) is affine in \(w\). The mean-ratio and weighted-mean-ratio constraints are linear in \(w\). Moreover, the PRD box constraints can be rearranged into linear inequalities on each fold. Indeed, the lower and upper PRD bounds are equivalent to
\begin{align}
    \label{eq:prd_linearized_bounds}
    \frac{1}{n_v}\sum_{i\in V_v}\frac{\hat P_i^{\,v}(w)}{P_i^{\,v}}
    -
    \underline{\tau}_{\mathrm{PRD}}
    \frac{\sum_{i\in V_v}\hat P_i^{\,v}(w)}{\sum_{i\in V_v}P_i^{\,v}}
    &\ge 0,
    \qquad \forall v\in\mathcal V,
    \\
    \frac{1}{n_v}\sum_{i\in V_v}\frac{\hat P_i^{\,v}(w)}{P_i^{\,v}}
    -
    \overline{\tau}_{\mathrm{PRD}}
    \frac{\sum_{i\in V_v}\hat P_i^{\,v}(w)}{\sum_{i\in V_v}P_i^{\,v}}
    &\le 0,
    \qquad \forall v\in\mathcal V,
    \nonumber
\end{align}
and both sides are affine in \(w\). Consequently, \eqref{opt:stacking_convex_example} is a convex quadratic program over the simplex.

This example illustrates the main appeal of the stacking extension. In the single-model case, calibration sorts a finite set of candidate configurations and selects the best feasible one. In the stacking case, the same assessor-facing diagnostics can be enforced while optimizing continuously over nonnegative combinations of candidate predictors, which can enlarge the attainable accuracy--equity trade-off frontier.

% \paragraph{Remark on COD.}
% Exact COD is not included in \eqref{opt:stacking_convex_example} because it depends on the foldwise median ratio, and that median dependence generally breaks the clean convexity structure above. For this reason, COD fits naturally into the discrete single-model screening problem \eqref{opt:model_selection_single}, whereas a convex stacking formulation is most transparent when built from diagnostics such as mean ratio, weighted mean ratio, and PRD. If one wishes to include a dispersion requirement in the stacking problem while preserving convexity, a convex surrogate or alternative dispersion measure can be used instead.



\section{Constrained Model Selection Across Validation Folds}
\label{sec:model_selection}

The covariance-regularized learners developed in the previous sections enlarge the set of candidate specifications available at calibration time, but they do not change the basic structure of the assessor's workflow. After training the implemented model families---namely regularized linear models and gradient boosting models, with or without covariance-based penalization---the practical task remains to compare a finite collection of candidate specifications on out-of-sample validation data and to select one final specification for deployment. Because the resulting tradeoffs involve predictive performance, ratio-study diagnostics, regressivity measures, and dispersion measures, we formulate this step as a constrained single-model selection problem.

Let $\mathcal V=\{1,\dots,V\}$ index the validation folds (or rolling-origin validation windows), and let $V_v$ denote the validation sample associated with fold $v\in\mathcal V$. Let $\mathcal C$ denote a finite set of candidate model configurations. Each configuration $c\in\mathcal C$ specifies
(i) a model family from the implemented classes,
(ii) a penalty choice
\[
\Psi_n \in \{0,\Psi_n^{\mathrm{cov}^2},\Psi_n^{\mathrm{neg\mbox{-}cov}},\Psi_n^{\mathrm{surr}}\},
\]
where $\Psi_n=0$ denotes the unpenalized baseline, and
(iii) all associated tuning parameters, including the penalty level $\rho\ge 0$ when applicable.
For each fold $v\in\mathcal V$, fitting configuration $c$ on the corresponding training portion yields a predictor $f_{c,v}$.

For each observation $i\in V_v$, let
\[
\hat y_i^{\,v}(c):=f_{c,v}(x_i^{\,v}),
\qquad
\hat P_i^{\,v}(c):=\exp\!\bigl(\hat y_i^{\,v}(c)\bigr),
\qquad
r_i^{\,v}(c):=\frac{\hat P_i^{\,v}(c)}{P_i^{\,v}},
\]
where $P_i^{\,v}$ is the observed sale price. From the validation predictions on fold $v$, we compute a predictive performance score $\mathrm{Perf}_v(c)$, written so that larger values are preferred, together with a collection of fold-level assessment diagnostics $Q_{m,v}(c)$ indexed by $m$. If one starts from an error metric $E_v(c)$ that is minimized, such as MSE, RMSE, or MAE, one may simply set $\mathrm{Perf}_v(c):=-E_v(c)$.

To aggregate results across folds, let $\alpha_v\ge 0$ satisfy $\sum_{v\in\mathcal V}\alpha_v=1$, and define
\[
\bar{\mathrm{Perf}}(c):=\sum_{v\in\mathcal V}\alpha_v\,\mathrm{Perf}_v(c),
\qquad
\bar Q_m(c):=\sum_{v\in\mathcal V}\alpha_v\,Q_{m,v}(c).
\]
The default choice is uniform weighting, $\alpha_v=1/V$, although the same formulation also allows non-uniform weights when one wishes to emphasize, for example, more recent validation windows.

We then select a single deployed model by solving
\begin{align}
    \label{opt:model_selection}
    c^\star
    \in
    \arg\max_{c\in\mathcal C}
    \quad
    & \bar{\mathrm{Perf}}(c)
    \\
    \text{s.t.}\quad
    & \underline{\kappa}_m \le \bar Q_m(c) \le \overline{\kappa}_m,
    \qquad m\in\mathcal M_{\mathrm{int}},
    \nonumber\\
    & \bar Q_m(c) \le \overline{\kappa}_m,
    \qquad m\in\mathcal M_{\mathrm{up}}.
    \nonumber
\end{align}
Here, $\mathcal M_{\mathrm{int}}$ indexes diagnostics constrained to lie in admissible intervals, while $\mathcal M_{\mathrm{up}}$ indexes diagnostics handled through upper bounds only. Typical elements of $\mathcal M_{\mathrm{int}}$ include ratio statistics such as the median ratio, mean ratio, and weighted mean ratio, as well as regressivity diagnostics such as VEI, PRD, and PRB. Typical elements of $\mathcal M_{\mathrm{up}}$ include dispersion measures such as COD and COV. Additional assessor-facing diagnostics can be incorporated by adding them to either set according to the desired feasibility requirement.

Formulation \eqref{opt:model_selection} deliberately avoids collapsing all validation criteria into a single composite score. Predictive performance remains the objective, while ratio-study, regressivity, and dispersion diagnostics enter as explicit constraints. This mirrors the logic of assessor calibration more closely than a weighted-score approach: among all candidate specifications that satisfy the chosen assessment-quality requirements, select the one with the best aggregated out-of-sample predictive performance.

The formulation above uses aggregated fold-level diagnostics as the default choice. When temporal stability is especially important, one may replace some aggregated constraints by foldwise requirements, or use a more conservative aggregation rule. We keep the weighted-average formulation as the main presentation because it yields a transparent calibration problem over a finite candidate set and aligns naturally with standard validation practice.

Once an optimal configuration $c^\star$ has been identified, the final deployed predictor is obtained by refitting that same configuration on the full pre-assessment development sample. In this sense, the procedure remains a single-model selection step: it chooses one specification from the implemented model families and one associated set of tuning parameters, rather than forming a composite predictor.

A natural extension, which we do not pursue here, is to replace the discrete feasible set $\mathcal C$ by convex combinations of candidate prediction vectors. Such a convex-stacking extension would preserve the same objective--constraint logic while potentially enlarging the attainable accuracy--equity trade-off frontier.





\section{Empirical Study: Cook County Assessor's Office}
\label{sec:empirical}

\subsection{Scope, Pipeline, and Data}

The empirical study is based on the Cook County Assessor's Office workflow (CCAO). The CCAO manages extensive documentation and a Github repository \todo{add Github link as a citation}. For the experiments, we follow the documentation, available in the \texttt{model-res-avm} codebase and, specifically, the version available for 2025 assessment process (\texttt{assessment-year-2025}). This includes data preparation, preprocessing, and feature engineering, mostly implemented in the R programming language. However, our implementation is in Python~3.10.14 \todo{cite?}, as a translation of the original R-based workflow. The temporal splits are chosen to mimic a realistic assessor setting: 2016--2022 form the train/validation development period, 2023 is used as the out-of-time test year, and 2024 is treated as the current assessment year reserved for later deployment analysis (see for specific details  \todo{Is this necessary? I was working with 2023 as test and 2024 as assessment, but not sure if makes sense for the paper}). 

\todo{
At the current stage of the project, the empirical evidence reported in the manuscript comes from the validation and test components of this translated pipeline. 

The final paper should continue to state clearly whether any remaining differences persist between the translated implementation and the original production workflow, since that distinction matters for both reproducibility and external interpretation.}


% \begin{table}[!ht]
% \centering
% \caption{Temporal data splits in the CCAO-derived pipeline.}
% \label{tab:data_splits}
% \begin{tabular}{lll}
% \toprule
% Split & Begin date & End date \\
% \midrule
% Train/Validation & Jan.\ 01, 2016 & Dec.\ 31, 2022 \\
% Test & Jan.\ 01, 2023 & Dec.\ 31, 2023 \\
% Assessment & Jan.\ 01, 2024 & Dec.\ 27, 2024 \\
% \bottomrule
% \end{tabular}
% \end{table}

\subsection{Preprocessing, Filtering, and Sample Construction}

\todo{
Because the empirical contribution is intended as a real workflow study rather than a purely conceptual example, the final paper should report the main inclusion rules, exclusion rules, and final sample sizes used in the translated CCAO pipeline. Table~\ref{tab:sample_construction_placeholder} is a placeholder for that information. At a minimum, the final version should document the property scope, any multicard or outlier exclusions, the construction of the sale-price target, and the final observation counts by temporal split.
}


\begin{table}[!ht]
\centering
\caption{Placeholder for sample construction summary. Replace with the final counts and main filtering rules from the experiment pipeline.}
\label{tab:sample_construction_placeholder}
\fbox{\parbox{0.9\textwidth}{
\textbf{Insert final sample-construction table here.} Suggested columns: stage, main rule applied, number of remaining observations, and split-specific counts. Suggested rows: raw extract, sale verification filters, multicard exclusion, outlier exclusion, modeling-ready sample, and split counts for train/validation, test, and assessment.
}}
\end{table}

\subsection{Exact Validation Windows and Fold Weights}

\todo{
The constrained calibration framework in Section~\ref{sec:model_selection} is defined over validation folds, so the final paper should report the exact rolling-origin or expanding-window construction used inside the 2016--2022 development period. Table~\ref{tab:fold_definition_placeholder} is a placeholder for that exact specification. This should include the training window, validation window, whether the training sample expands over time, and the fold weights \(\alpha_v\) used in metric aggregation.
}

\begin{table}[!ht]
\centering
\caption{Placeholder for exact validation-window construction.}
\label{tab:fold_definition_placeholder}
\fbox{\parbox{0.9\textwidth}{
\textbf{Insert final fold-definition table here.} Suggested columns: fold index, training years, validation years, number of validation observations, and aggregation weight \(\alpha_v\). This table should make the calibration protocol fully reproducible.
}}
\end{table}

\subsection{Baselines: Linear Versus Nonlinear Models}


As mentioned in Section \todo{Check}, the CCAO used Linear Regression-based prediction models for the mass appraisal process until 2019 \todo{check}, and then transitioned to the nonlinear boosted tree model LightGBM. In terms of accuracy, the transition is consistent with the literature; boosted trees tend to perform better on tabular data \todo{cite}. 

Tables~\ref{tab:baseline_models} and~\ref{tab:baseline_models_2} report the current benchmark comparison for the CCAO's data on the test set and assessment sets. The baseline LightGBM model substantially improves predictive accuracy relative to linear regression. However, at the same time it worsens vertical-equity diagnostics. Relative to Linear Regression, the baseline LightGBM model improves \(R^2\), MAE, COD, and COV, but its VEI, PRD, and PRB values are worsening, until the point they are clearly signaling a regressive behavior, while Linear regression seems much more stable in this dimension.

\medskip

\noindent \textbf{Remark.} This is the main empirical motivation for the paper: stronger nonlinear predictive fit does not automatically imply better vertical equity. This not a surprising behavior; nonlinear models can capture much more stronger relationships, and thus are more susceptible to produce non-equitable patters like in (\label{eq:covariance_population_identity}) \todo{Verify it?}\todo{Cite maybe?}. Therefore, there is a necessity for a middle-ground solution between a simple but equitable model and more complex one that can achieve higher prediction performance, which motivates our solution.
% . On the other hand, while Linear Regression worsens the accuracy and dispersion, it seems to satisfy all the desired values for the regressivity measures  introduced in Table \ref{tab:assessment_metrics_summary}. 

% \begin{table}[t]
% \centering
% \caption{Baseline comparison on the test set. Lower is better for MAE, COD, and COV. For VEI and PRB, values closer to zero are preferred; for PRD and MKI, values closer to one are preferred.}
% \label{tab:baseline_models}
% \resizebox{\textwidth}{!}{%
% \begin{tabular}{lrrrrrrrr}
% \toprule
% Model & \(R^2\) & MAE & COD & COV & VEI & PRD & PRB & MKI \\
% \midrule
% Linear Regression & 0.798 & \$85{,}273 & 24.306 & 0.354 & 6.223\% & 1.016 & 0.033 & 1.009 \\
% LGBM Regressor & 0.860 & \$73{,}398 & 21.647 & 0.330 & -31.431\% & 1.082 & -0.108 & 0.890 \\
% \bottomrule
% \end{tabular}}
% \end{table}

\begin{table}[!ht]
\centering
\caption{Baseline comparison on the \textit{test} set ($\sim$2023). Lower is better for MAE, COD, and COV. For PRD, values closer to one are preferred; for VEI and PRB, values closer to zero are preferred.}
\label{tab:baseline_models}
\resizebox{\textwidth}{!}{%
\begin{tabular}{lrrrrrrr}
\toprule
Model & \(R^2\) & MAE & COD & COV & PRD & PRB & VEI \\
\midrule
Linear Regression & 0.762 & \$86{,}257 & 25.013\% & 36.825\% & 1.039 & -0.011 & -13.474\% \\
LightGBM & 0.887 & \$71{,}337 & 22.739\% & 34.805\% & 1.082 & -0.114 & -36.810\% \\
% LGBCovPenalty_mode_diff_rho_1.0 & 0.886 & \$72{,}286 & 23.018\% & 35.083\% & 1.076 & -0.105 & -32.211\% \\
% LGBSmoothPenalty_mode_diff_rho_1.0 & 0.888 & \$72{,}838 & 22.568\% & 33.896\% & 1.063 & -0.069 & -21.310\% \\
\bottomrule
\end{tabular}}
\end{table}



\begin{table}[!ht]
\centering
\caption{Baseline comparison on the \textit{assessment} set (2024). Lower is better for MAE, COD, and COV. For PRD, values closer to one are preferred; for VEI and PRB, values closer to zero are preferred.}
\label{tab:baseline_models_2}
\resizebox{\textwidth}{!}{%
\begin{tabular}{lrrrrrrr}
\toprule
Model & \(R^2\) & MAE & COD & COV & PRD & PRB & VEI \\
\midrule
LinearRegression & 0.802 & \$90{,}749 & 24.199\% & 36.622\% & 1.042 & -0.022 & -13.352\% \\
LGBMRegressor & 0.886 & \$79{,}816 & 21.310\% & 35.097\% & 1.071 & -0.096 & -28.208\% \\
% LGBCovPenalty_mode_diff_rho_1.0 & 0.888 & \$80{,}331 & 21.515\% & 35.408\% & 1.062 & -0.081 & -21.184\% \\
% LGBSmoothPenalty_mode_diff_rho_1.0 & 0.885 & \$81{,}863 & 21.252\% & 34.661\% & 1.048 & -0.046 & -11.638\% \\
\bottomrule
\end{tabular}}
\end{table}


\todo{Table~\ref{tab:baseline_models} should be paired with the main motivation plot in the final version. Figure~\ref{fig:baseline_regressivity_placeholder} is the placeholder for that visual.
Why would we add the plot?}


\begin{figure}[t]
\centering
\fbox{\parbox{0.9\textwidth}{
\textbf{Insert baseline regressivity plot here.} Suggested figure: assessment ratio \(r=\hat P/P\) versus log market value under the baseline LightGBM model, including the horizontal \(r=1\) line, a LOWESS trend, and an optional linear trend line. This is the main visual motivation for the regressivity problem.
}}
\caption{Placeholder for the baseline regressivity motivation figure.}
\label{fig:baseline_regressivity_placeholder}
\end{figure}

\subsection{Effect of Covariance Penalization}

We aim to get a middle-ground solution that can achieve better tradeoffs between accuracy and vertical-equity, when compared to the baseline models (Linear Regression and LightGBM). We propose to add covariance-based penalties to the baseline LightGBM learner. The current empirical results show the basic mechanism clearly \todo{CHECK THE NUMBERS}. For the baseline model, the ratio--log-value relationship exhibits a strong negative trend, with a reported linear slope of approximately \(-0.2156\) and ratio--price correlation of approximately \(-0.2666\). Under the squared covariance penalty, the slope moves to about \(-0.1263\) and the correlation to about \(-0.1224\). Under the separable surrogate penalty, the slope moves even closer to flat, about \(-0.0222\), and the correlation to about \(-0.0239\). These are not deployment diagnostics by themselves, but they visually illustrate that the training-time penalties act in the intended direction.


The final paper should show that before/after comparison directly. Figure~\ref{fig:penalty_mechanism_placeholder} is the placeholder for the main paired ratio-trend figure.



\begin{figure}[!htbp]
    \centering
    % First subfigure
    \begin{subfigure}{0.48\textwidth}
        \centering
        \includegraphics[width=\linewidth, trim={-3cm 0mm 3cm 3mm}]{Fairness Price/CCAO meeting/img/motivation/Baseline LightGBM_0_motivation.pdf}
        \caption{\scriptsize Baseline LightBGM}
    \end{subfigure}\hfill
    % Second subfigure
    \begin{subfigure}{0.48\textwidth}
        \centering
        \includegraphics[width=\linewidth, trim={3cm 0mm -3cm 3mm}]{Fairness Price/CCAO meeting/img/Covariance-Penalized LightGBM_2_motivation.pdf}
        \caption{\scriptsize Covariance-Penalized LightGBM}
    \end{subfigure}
    
    \caption{LightGBM example comparison. / Placeholder for the main mechanism figure showing how covariance penalization flattens the ratio--value trend. /  Suggested layout: two panels comparing baseline LightGBM against a squared-covariance or surrogate-penalized LightGBM, each showing ratio versus log value, the \(r=1\) line, and trend overlays. Add the reported slope and correlation summaries in the caption or callout box.}
\end{figure}
\vspace{-5mm}

% \begin{figure}[t]
% \centering
% \fbox{\parbox{0.9\textwidth}{
% \textbf{Insert before/after penalty mechanism figure here.} Suggested layout: two panels comparing baseline LightGBM against a squared-covariance or surrogate-penalized LightGBM, each showing ratio versus log value, the \(r=1\) line, and trend overlays. Add the reported slope and correlation summaries in the caption or callout box.
% }}
% \caption{Placeholder for the main mechanism figure showing how covariance penalization flattens the ratio--value trend.}
% \label{fig:penalty_mechanism_placeholder}
% \end{figure}

\subsection{Penalized LightGBM Performance on the Test Set}

Table~\ref{tab:penalized_models} summarizes the current test-set results for the main penalized LightGBM variants. All four reported penalized models move the vertical-equity diagnostics sharply toward assessor-preferred ranges. For example, the baseline LightGBM has PRD \(=1.082\) and PRB \(=-0.108\), both strongly regressive, whereas the penalized variants move PRD into the 1.029--1.030 range and PRB into the \(-0.026\) to \(0.008\) range. At the same time, the best reported squared-covariance configuration improves \(R^2\) from \(0.860\) to \(0.880\) and improves MAE from \$73{,}398 to \$70{,}855. These results are exactly the kind of middle-ground behavior the method is designed to find.

% \begin{table}[t]
% \centering
% \caption{Covariance-penalized LightGBM results on the test set. The reported \(\rho\) values are penalty-specific tuning parameters and are not directly comparable across penalty definitions.}
% \label{tab:penalized_models}
% \resizebox{\textwidth}{!}{%

% \begin{tabular}{llrrrrrrr}
% \toprule
%  & Model & R$^2$ & MAE & COD & COV & VEI & PRD & PRB \\
% \midrule
% & Linear Regression & 0.802044 & 90749.426470 & 24.199356 & 0.366218 & -13.351674 & 1.042389 & -0.021762 \\
%  & LightGBM & 0.885557 & 79815.915806 & 21.310066 & 0.350968 & -28.207672 & 1.070587 & -0.095780 \\
% % 2 & LGBCovPenalty_mode_diff_rho_1.0 & 0.887535 & 80331.081744 & 21.514531 & 0.354079 & -21.183510 & 1.061787 & -0.081160 \\
% % 3 & LGBSmoothPenalty_mode_diff_rho_1.0 & 0.885457 & 81863.289919 & 21.251802 & 0.346613 & -11.637937 & 1.048117 & -0.045758 \\
% \bottomrule
% \end{tabular}

% % \begin{tabular}{lrrrrrrrrr}
% % \toprule
% % Model & \(\rho\) & \(R^2\) & MAE & COD & COV & VEI & PRD & PRB & MKI \\
% % \midrule
% % LGBM Regressor & -- & 0.860 & \$73{,}398 & 21.647 & 0.330 & -31.431\% & 1.082 & -0.108 & 0.890 \\
% % LGBMSurrPenalty [v1] & 2.043 & 0.872 & \$74{,}579 & 21.770 & 0.316 & -0.043\% & 1.030 & 0.003 & 0.970 \\
% % LGBMSurrPenalty [v2] & 280.7 & 0.867 & \$75{,}641 & 21.829 & 0.313 & 3.403\% & 1.029 & 0.008 & 0.970 \\
% % LGBMCovPenalty [v1] & 6.58 & 0.877 & \$72{,}042 & 21.693 & 0.319 & -6.579\% & 1.030 & -0.026 & 0.989 \\
% % LGBMCovPenalty [v2] & 1487.0 & 0.880 & \$70{,}855 & 21.467 & 0.317 & -4.043\% & 1.029 & -0.022 & 0.986 \\
% % \bottomrule
% % \end{tabular}
% }
% \end{table}




\begin{table}[htbp]
    \centering
    \caption{\textbf{Penalized Objective Comparison:} Pred. Error + $\rho\; \cdot$ Covariance(ratios, prices)$^2$. Test results for Accuracy (R2, MAE), horizontal equity (COD, COV), and vertical equity (VEI, PRD, PRB, MKI). / / Covariance-penalized LightGBM results on the test set. The reported \(\rho\) values are penalty-specific tuning parameters and are not directly comparable across penalty definitions.}
    \label{tab:penalized_models}
    \resizebox{\textwidth}{!}{%
    \begin{tabular}{lc|cc|cc|cccc}
    \toprule
     Model & Penalty $\rho$ & R2 & MAE & COD & COV (IAAO) & VEI & PRD & PRB & MKI \\
    \midrule
    % --- Baselines ---
    LinearRegression & -- & 0.798 & \$85,273 & 24.306 & 0.354 & 6.223\% & 1.016 & 0.033 & 1.009 \\
    LGBMRegressor & -- & 0.860 & \$73,398 & 21.647 & 0.330 & \textcolor{red}{-31.431\%} & \textcolor{red}{1.082} & \textcolor{red}{-0.108} & \textcolor{red}{0.890} \\
    \midrule
    % --- Penalized Models ---
    LGBMSurrPenalty [v1] & 2.043 &
    \makecell{\textbf{0.872}\\{\scriptsize(\textcolor{green!60!black}{+1.40\%})}} &
    \makecell{\$74,579\\{\scriptsize(\textcolor{red!70!black}{-1.61\%})}} &
    \makecell{21.770\\{\scriptsize(\textcolor{red!70!black}{-0.57\%})}} &
    \makecell{\textbf{0.316}\\{\scriptsize(\textcolor{green!60!black}{+4.24\%})}} &
    \makecell{\textbf{-0.043\%}\\{\scriptsize(\textcolor{green!60!black}{+99.86\%})}} &
    \makecell{\textbf{1.030}\\{\scriptsize(\textcolor{green!60!black}{+63.41\%})}} &
    \makecell{\textbf{+0.003}\\{\scriptsize(\textcolor{green!60!black}{+97.22\%})}} &
    \makecell{\textbf{0.970}\\{\scriptsize(\textcolor{green!60!black}{+72.73\%})}} \\
    [1.5mm]
    
    LGBMSurrPenalty [v2] & 280.7 &
    \makecell{\textbf{0.867}\\{\scriptsize(\textcolor{green!60!black}{+0.81\%})}} &
    \makecell{\$75,641\\{\scriptsize(\textcolor{red!70!black}{-3.06\%})}} &
    \makecell{21.829\\{\scriptsize(\textcolor{red!70!black}{-0.84\%})}} &
    \makecell{\textbf{0.313}\\{\scriptsize(\textcolor{green!60!black}{+5.15\%})}} &
    \makecell{\textbf{+3.403\%}\\{\scriptsize(\textcolor{green!60!black}{+89.17\%})}} &
    \makecell{\textbf{1.029}\\{\scriptsize(\textcolor{green!60!black}{+64.63\%})}} &
    \makecell{\textbf{+0.008}\\{\scriptsize(\textcolor{green!60!black}{+92.59\%})}} &
    \makecell{\textbf{0.970}\\{\scriptsize(\textcolor{green!60!black}{+72.73\%})}} \\
    [1.5mm]
    
    LGBMCovPenalty [v1] & 6.58 &
    \makecell{\textbf{0.877}\\{\scriptsize(\textcolor{green!60!black}{+1.98\%})}} &
    \makecell{\textbf{\$72,042}\\{\scriptsize(\textcolor{green!60!black}{+1.85\%})}} &
    \makecell{21.693\\{\scriptsize(\textcolor{red!70!black}{-0.21\%})}} &
    \makecell{\textbf{0.319}\\{\scriptsize(\textcolor{green!60!black}{+3.33\%})}} &
    \makecell{\textbf{-6.579\%}\\{\scriptsize(\textcolor{green!60!black}{+79.07\%})}} &
    \makecell{\textbf{1.030}\\{\scriptsize(\textcolor{green!60!black}{+63.41\%})}} &
    \makecell{\textbf{-0.026}\\{\scriptsize(\textcolor{green!60!black}{+75.93\%})}} &
    \makecell{\textbf{0.989}\\{\scriptsize(\textcolor{green!60!black}{+90.00\%})}} \\
    [1.5mm]
    
    LGBMCovPenalty [v2] & 1487.0 &
    \makecell{\textbf{0.880}\\{\scriptsize(\textcolor{green!60!black}{+2.33\%})}} &
    \makecell{\textbf{\$70,855}\\{\scriptsize(\textcolor{green!60!black}{+3.46\%})}} &
    \makecell{\textbf{21.467}\\{\scriptsize(\textcolor{green!60!black}{+0.83\%})}} &
    \makecell{\textbf{0.317}\\{\scriptsize(\textcolor{green!60!black}{+3.94\%})}} &
    \makecell{\textbf{-4.043\%}\\{\scriptsize(\textcolor{green!60!black}{+87.14\%})}} &
    \makecell{\textbf{1.029}\\{\scriptsize(\textcolor{green!60!black}{+64.63\%})}} &
    \makecell{\textbf{-0.022}\\{\scriptsize(\textcolor{green!60!black}{+79.63\%})}} &
    \makecell{\textbf{0.986}\\{\scriptsize(\textcolor{green!60!black}{+87.27\%})}} \\
    \bottomrule
    \end{tabular}
    }
\end{table}



If desired, the final paper can also add a second table with percentage improvements relative to the baseline LightGBM model. That presentation is useful for quick interpretation, but Table~\ref{tab:penalized_models} should remain the primary source of the raw metric values.

\subsection{Tradeoff Frontiers}

The current results indicate that increasing \(\rho\) traces a meaningful accuracy--equity tradeoff curve rather than producing a degenerate monotone sacrifice of fit. In PRD--\(R^2\) space, the penalized models move from the unconstrained LightGBM point toward the assessor compliance band while remaining near the upper edge of the accuracy frontier. A similar pattern appears in VEI--\(R^2\) space. These frontier views justify the constrained calibration logic of Section~\ref{sec:model_selection}: in a genuinely multi-criteria setting, the practical question is not which model is globally best, but which feasible model offers the best accuracy among those that satisfy the chosen assessment-quality requirements.

\begin{figure}[!htbp]
    \centering
    % First Subfigure
    \begin{subfigure}[b]{0.48\textwidth}
        \centering
        \includegraphics[width=\textwidth, trim={2mm 0 4mm 2mm}]{Fairness Price/CCAO meeting/img/tradeoffs/PRD_vs_R2_band.pdf}
        \caption{Insert PRD description here}
        \label{fig:prd_tradeoff}
    \end{subfigure}
    \hfill
    % Second Subfigure
    \begin{subfigure}[b]{0.48\textwidth}
        \centering
        \includegraphics[width=\textwidth, trim={4mm 0 0 2mm}]{Fairness Price/CCAO meeting/img/tradeoffs/VEI_vs_R2_band.pdf}
        \caption{Insert VEI description here}
        \label{fig:vei_tradeoff}
    \end{subfigure}
    
    % Main Figure Caption and Label
    \caption{\textbf{Pareto Frontier:} Test set behavior for different penalization terms ($\rho >0$).}
    \label{fig:pareto_frontier}
\end{figure}

\begin{figure}[!ht]
\centering
\includegraphics[width=0.7\linewidth, trim={2mm 2mm 4mm 4mm}]{Fairness Price/CCAO meeting/img/model_selection/pareto_optimal.pdf}
\caption{\parbox{0.9\textwidth}{ Suggested layout: one panel for PRD versus \(R^2\) with the compliance band indicated, and one panel for VEI versus \(R^2\). Label baseline linear regression, baseline LightGBM, and the penalized families across increasing \(\rho\).
}}
% \caption{Placeholder for the main accuracy--equity frontier figure.}
\label{fig:frontier_placeholder}
\end{figure}

% \begin{figure}
%         \caption{{\textbf{Pareto Frontier:} Test set behavior for different penalization terms ($\rho >0$)}}
%         \centering
        
% \end{figure}

\subsection{Validation-Fold Constrained Calibration}

One of the paper's methodological contributions is the constrained calibration rule of Section~\ref{sec:model_selection}, so the empirical section should also instantiate that rule directly. The final version should report the exact active validation constraints, the number of candidate configurations evaluated, the number of feasible configurations, the selected configuration, and the corresponding performance on the untouched test year. Table~\ref{tab:calibration_constraints_placeholder} and Table~\ref{tab:selected_model_placeholder} are placeholders for that information.

\begin{table}[t]
\centering
\caption{Placeholder for the exact empirical calibration rule used on validation folds.}
\label{tab:calibration_constraints_placeholder}
\fbox{\parbox{0.9\textwidth}{
\textbf{Insert final calibration-constraints table here.} Suggested columns: diagnostic, constraint type, lower bound, upper bound, and whether the constraint is aggregated across folds or enforced foldwise. Typical rows may include PRD, PRB, VEI, COD, COV, and ratio-statistic guardrails.
}}
\end{table}

\begin{table}[t]
\centering
\caption{Placeholder for the final selected configuration under constrained calibration.}
\label{tab:selected_model_placeholder}
\fbox{\parbox{0.9\textwidth}{
\textbf{Insert selected-model summary here.} Suggested rows or columns: selected model family, penalty type, \(\rho\), number of feasible candidates, best infeasible accuracy-only model, final selected feasible model, and held-out test metrics for the selected configuration.
}}
\end{table}

This subsection is also the natural place to explain any difference between (i) the full frontier of candidate models and (ii) the single final model selected for deployment under the chosen assessment-quality requirements.

\subsection{Uncertainty, Stability, and Fold Heterogeneity}

Because regressivity diagnostics can be noisy across time windows, the final paper should also report a stability summary rather than relying only on point estimates. This can be done through bootstrap confidence intervals on the held-out test metrics, foldwise distributions across validation windows, or both. Figure~\ref{fig:stability_placeholder} is the placeholder for that summary.

\begin{figure}[t]
\centering
\fbox{\parbox{0.9\textwidth}{
\textbf{Insert stability/uncertainty figure here.} Suggested options: foldwise boxplots for PRD, PRB, VEI, COD, and \(R^2\); bootstrap confidence intervals for the selected model versus the main baselines; or a table of foldwise means and standard deviations.
}}
\caption{Placeholder for stability and uncertainty reporting.}
\label{fig:stability_placeholder}
\end{figure}

\subsection{Practical Implications for Assessor Workflows}

The empirical message is operationally simple. The method can be inserted into an existing valuation workflow with minimal changes. It uses the same preprocessing, the same base model classes, and the same downstream diagnostics already familiar to assessors. It adds one main penalty parameter with a transparent qualitative effect: larger \(\rho\) values tend to move models toward less regressive behavior. This makes the method attractive not only because it improves the attainable tradeoffs, but because it does so in a form that can be calibrated, audited, and communicated within existing assessor practice.





% OLD TEXT


\section{(Optional) important points of the approach. Other have empirical examination or case of study can be an option.}


% \section{Experiments Results}




% \section{Extensions}

\section{Discussion/Conclusions, Limitations, Extensions \& Future Work}

\subsection{Future Work}

Same direction: control the relationship, but now aware of the sub-market clusters, which can create nonlinear relationships when asking for too much overall linear correction (see image of 0-correlation plot of surrogate penalty).

Interpretable \& fair models. Apply this to multiple cases of assessors offices, neural networks as n alternative methodology because of their predictive power (?); in particular, their ability to handle Geospatial data (e.g., apply an embedding using the information or methodology provided by AlphaEarth of Google). Another alternative is to test the foundational models (e.g., tabular models). Try to incorpore it into Random Forest/Bagging, which is also widely used apparently.

\section{MY NOTES}

Paper structure references. INFORMS seems to use a lot of plots and tables in general (make sense since they care about the analytics part).
\begin{itemize}
    \item INFORMS data science: https://pubsonline.informs.org/doi/epdf/10.1287/ijds.2024.0029
    \item INFORMS managements science: https://pubsonline.informs.org/doi/epdf/10.1287/mnsc.2022.00345
    \item INFORM journal on optimization: https://arxiv.org/pdf/2412.17116
\end{itemize}



\noindent Haihao whiteboard:
\begin{itemize}
    \item Introduction: motivation, contribution, related work.
    \item Setting: Examples, assumptions, etc.
    \item Methodology
    \item Experiments 
    \item Extension (2-pages at most)
    \item Conclusions
\end{itemize}

\noindent Swati Comments / Suggestions:

\begin{itemize}
    \item What is novel in this research (theory and applications-wise)? Has anything like this be done already?
    \item Good thing: it looked like I owned the project. Lets keep it that way for the GE.
    \item It has to be a journal-level paper. Good venues: Journal of Machine Learning Research (JMLR) or Management Sciences. 
    \item Paper Critique: adjacent to what I've done. Maybe ML, fairness, etc. The idea is to be able to read, summarize, critique (Not just chatGPT-style). Even go on details, proofs, technical stuff. I can do extra research on the topic of the paper.
    \item Important: GE as a message of what I've been doing during the phd.
    \item Presentation (oral-exam): Be prepared for presenting + questions at the same time. Bad things about past presentations: too much notation, not be able to answer extra questions, etc.
    \item Benchmark v/s more complex correlation/covariance methods. They don't work/generalize as good. Show it.
    \item Extras: copy of the paper \& slides in physical for reading. 
\end{itemize}

\printbibliography %Prints bibliography


% APPENDIX
\newpage

\appendix


\section{A Theoretical Analysis of Regressivity as a Dependence Problem}
\label{anex:regressivity_dependence}

The assessor literature describes \emph{regressivity} as a price-related pattern in which assessment ratios tend to decline as property value increases \cite{IAAO2013Ratio,IAAO2025ExposureRatio}. In the predictive setting of Section~\ref{subsec:prediction_setting}, this pattern can be expressed directly in terms of the model outputs.

For a fitted model $f$ and an evaluation sample $S$, recall the assessment ratio
\[
r_i(f):=\frac{\widehat P_i(f)}{P_i},
\qquad i\in S,
\]
and define the log-ratio residual
\[
u_i(f):=\log r_i(f)=\hat y_i(f)-y_i,
\]
where the identity follows from \eqref{eq:log_ratio_identity}. Thus, any price-related pattern in the assessment ratios can equivalently be studied as a dependence pattern between $u_i(f)$ and a value index.

To keep the exposition slightly more general, let
\[
z_i:=\phi(P_i),
\]
where $\phi:\mathbb{R}_{++}\to\mathbb{R}$ is any strictly increasing transform of price. Throughout the paper the default choice is $\phi(P)=\log P$, but the arguments below only require monotonicity. For any pair of sample vectors $a=\{a_i\}_{i\in S}$ and $b=\{b_i\}_{i\in S}$, define
\begin{equation}
\label{eq:sample_cov_general}
\widehat{\operatorname{Cov}}_S(a,b)
:=
\frac{1}{|S|}
\sum_{i\in S}
(a_i-\bar a_S)(b_i-\bar b_S),
\end{equation}
with the analogous definition for $\widehat{\operatorname{Var}}_S(b)$. A natural log-scale summary of directional regressivity is therefore
\begin{equation}
\label{eq:cov_log_resid}
\widehat{\operatorname{Cov}}_S(u(f),z),
\end{equation}
and the corresponding linear trend is the slope
\begin{equation}
\label{eq:sample_slope_logz}
\widehat\beta^{\,\log}_S(f)
:=
\frac{\widehat{\operatorname{Cov}}_S(u(f),z)}
{\widehat{\operatorname{Var}}_S(z)}.
\end{equation}
Because $\widehat{\operatorname{Var}}_S(z)$ is fixed once the evaluation sample is fixed, controlling the sign and magnitude of the linear price-related trend is equivalent to controlling the covariance in \eqref{eq:cov_log_resid}.

At the same time, implementation is often simpler on the original ratio scale around the target of perfect equity. For this reason, define the price-scale deviation
\[
e_i(f):=r_i(f)-1.
\]
Since subtracting a constant does not change covariance,
\[
\widehat{\operatorname{Cov}}_S(e(f),z)
=
\widehat{\operatorname{Cov}}_S(r(f),z).
\]
Moreover, because $\sum_{i\in S}(z_i-\bar z_S)=0$,
\begin{equation}
\label{eq:cov_linear_form}
\widehat{\operatorname{Cov}}_S(e(f),z)
=
\frac{1}{|S|}
\sum_{i\in S} e_i(f)\,(z_i-\bar z_S).
\end{equation}
This representation motivates a separable surrogate penalty.

\begin{lemma}[Exact surrogate upper bound]
\label{lem:surrogate_upper_bound}
Define
\begin{equation}
\label{eq:surrogate_penalty}
\Psi^{\mathrm{surr}}_S(f)
:=
\frac{1}{|S|}
\sum_{i\in S} e_i(f)^2\,(z_i-\bar z_S)^2.
\end{equation}
Then, for every fitted model $f$,
\begin{equation}
\label{eq:cov_surrogate_bound}
\widehat{\operatorname{Cov}}_S(e(f),z)^2
\le
\Psi^{\mathrm{surr}}_S(f).
\end{equation}
\end{lemma}

\begin{proof}
By \eqref{eq:cov_linear_form},
\[
\widehat{\operatorname{Cov}}_S(e(f),z)
=
\frac{1}{|S|}
\sum_{i\in S} a_i(f),
\qquad
a_i(f):=e_i(f)\,(z_i-\bar z_S).
\]
Hence, by Jensen's inequality for $x\mapsto x^2$,
\[
\widehat{\operatorname{Cov}}_S(e(f),z)^2
=
\left(
\frac{1}{|S|}
\sum_{i\in S} a_i(f)
\right)^2
\le
\frac{1}{|S|}
\sum_{i\in S} a_i(f)^2
=
\frac{1}{|S|}
\sum_{i\in S} e_i(f)^2\,(z_i-\bar z_S)^2,
\]
which is exactly \eqref{eq:cov_surrogate_bound}.
\end{proof}

Lemma~\ref{lem:surrogate_upper_bound} gives the optimization interpretation used later: $\Psi_S^{\mathrm{surr}}(f)$ is not merely heuristic, but an exact upper bound on the squared covariance between the price-scale deviation from perfect equity and the value index. In addition, unlike the covariance itself, the surrogate is separable across observations, which makes it more convenient for gradient-based learners and libraries that are built around per-observation derivatives.

To connect this dependence view to standard ratio-study metrics, we now formalize the common monotone picture of regressivity.

\begin{assumption}[Monotone price-related behavior]
\label{ass:monotone_regressivity}
For a fitted model $f$ and evaluation sample $S$, there exists a nonincreasing function
\[
g_f:\phi(\mathbb{R}_{++})\to\mathbb{R}_{++}
\]
such that
\[
r_i(f)=g_f(z_i),
\qquad i\in S.
\]
Equivalently, whenever $z_i\le z_j$, one has $r_i(f)\ge r_j(f)$.
\end{assumption}

Assumption~\ref{ass:monotone_regressivity} is the finite-sample counterpart of the usual assessor description that ratios tend to fall with value. Because both $\phi$ and $\log$ are increasing, it immediately implies that $u_i(f)=\log r_i(f)$ is also nonincreasing in $z_i$.

\begin{proposition}[Sign of covariances under monotone regressivity]
\label{prop:cov_signs}
Under Assumption~\ref{ass:monotone_regressivity},
\begin{align}
\widehat{\operatorname{Cov}}_S(r(f),z) &\le 0,
\label{eq:cov_r_z_negative}
\\
\widehat{\operatorname{Cov}}_S(u(f),z) &\le 0,
\label{eq:cov_u_z_negative}
\\
\widehat{\operatorname{Cov}}_S(r(f),P) &\le 0.
\label{eq:cov_r_p_negative}
\end{align}
If $g_f$ is not constant on $S$, then the inequalities are strict.
\end{proposition}

\begin{proof}
For any two sample vectors $a$ and $b$,
\begin{equation}
\label{eq:pairwise_cov_identity}
\widehat{\operatorname{Cov}}_S(a,b)
=
\frac{1}{2|S|^2}
\sum_{i,j\in S}
(a_i-a_j)(b_i-b_j).
\end{equation}
Applying \eqref{eq:pairwise_cov_identity} with $a_i=r_i(f)$ and $b_i=z_i$, each summand is nonpositive because $r_i(f)=g_f(z_i)$ with $g_f$ nonincreasing. This proves \eqref{eq:cov_r_z_negative}.

Next, since $u_i(f)=\log r_i(f)$ and $\log$ is increasing, the map $z\mapsto u(f)=\log(g_f(z))$ is also nonincreasing, so the same argument yields \eqref{eq:cov_u_z_negative}.

Finally, because $\phi$ is strictly increasing, its inverse preserves order on the sample, so $P_i$ is increasing in $z_i$. Therefore $r_i(f)$ is nonincreasing in $P_i$ as well, and another application of \eqref{eq:pairwise_cov_identity} gives \eqref{eq:cov_r_p_negative}.
\end{proof}

Proposition~\ref{prop:cov_signs} shows that the sign of the covariance in either log space or price space is aligned with the usual notion of regressivity. The next result connects this to the assessor-facing vertical-equity metrics of Section~\ref{subsec:equity_metrics}.

\begin{proposition}[Implications for ratio-study metrics]
\label{prop:metrics_alignment}
Assume Assumption~\ref{ass:monotone_regressivity}. Then the following hold.

\smallskip
\noindent
\emph{(i) Log-slope.} The slope of log ratios on the value index is nonpositive:
\[
\widehat\beta^{\,\log}_S(f)\le 0.
\]

\smallskip
\noindent
\emph{(ii) PRD.} Let $\bar r_S:=|S|^{-1}\sum_{i\in S} r_i(f)$ and $\bar P_S:=|S|^{-1}\sum_{i\in S} P_i$. Then
\begin{equation}
\label{eq:prd_cov_identity}
\operatorname{PRD}(S)
=
\frac{1}
{1+\widehat{\operatorname{Cov}}_S(r(f),P)/(\bar r_S\bar P_S)}.
\end{equation}
Consequently,
\[
\widehat{\operatorname{Cov}}_S(r(f),P)\le 0
\quad\Longrightarrow\quad
\operatorname{PRD}(S)\ge 1,
\]
with strict inequality whenever $\widehat{\operatorname{Cov}}_S(r(f),P)<0$.

\smallskip
\noindent
\emph{(iii) PRB.} Let $v_i$ denote the value proxy used in the PRB calculation, and assume that it preserves the ordering by value on $S$ (that is, $P_i\le P_j$ implies $v_i\le v_j$). Then
\[
\operatorname{PRB}(S)\le 0.
\]

\smallskip
\noindent
\emph{(iv) VEI.} Let the lowest-value and highest-value groups used in the VEI be formed by the same ordering proxy $v_i$. Then, under the sign convention of Section~\ref{subsec:equity_metrics},
\[
\operatorname{VEI}(S)\le 0.
\]
\end{proposition}

\begin{proof}
Part (i) follows immediately from \eqref{eq:sample_slope_logz}, since $\widehat{\operatorname{Var}}_S(z)>0$ on any nondegenerate evaluation sample and Proposition~\ref{prop:cov_signs} gives $\widehat{\operatorname{Cov}}_S(u(f),z)\le 0$.

For part (ii), note that
\[
\overline{r_W}
=
\frac{\sum_{i\in S} r_i(f)P_i}{\sum_{i\in S} P_i}
=
\frac{\overline{rP}_S}{\bar P_S},
\]
where $\overline{rP}_S:=|S|^{-1}\sum_{i\in S} r_i(f)P_i$. Since
\[
\overline{rP}_S
=
\bar r_S\bar P_S
+
\widehat{\operatorname{Cov}}_S(r(f),P),
\]
we obtain
\[
\operatorname{PRD}(S)
=
\frac{\bar r_S}{\overline{r_W}}
=
\frac{\bar r_S\bar P_S}
{\bar r_S\bar P_S+\widehat{\operatorname{Cov}}_S(r(f),P)},
\]
which is exactly \eqref{eq:prd_cov_identity}. The sign implication is immediate.

For part (iii), the PRB response variable $(r_i(f)-m_S)/m_S$ is an affine transform of $r_i(f)$ and therefore remains nonincreasing in $v_i$ whenever $r_i(f)$ is nonincreasing in value and the proxy $v_i$ preserves that ordering. The PRB slope is therefore nonpositive by the same covariance-slope argument used in part (i).

For part (iv), if the sample is sorted by an ordering proxy that preserves value ranking, then every observation in the highest-value group has ratio no larger than every observation in the lowest-value group under Assumption~\ref{ass:monotone_regressivity}. Hence the median ratio of the highest-value group is no larger than the median ratio of the lowest-value group, which implies $\operatorname{VEI}(S)\le 0$ under the definition in Section~\ref{subsec:equity_metrics}.
\end{proof}

Proposition~\ref{prop:metrics_alignment} should be read as a \emph{directional} result: under a monotone ratio-value relationship, the covariance target and the standard vertical-equity metrics all move in the same regressive or non-regressive direction. This is precisely the sense in which covariance acts as a root measure of regressivity in the present framework.

The next corollary describes the neutral case.

\begin{corollary}[Flattening to perfect equity]
\label{cor:neutralization}
Let $\{f_m\}_{m\ge 1}$ be a sequence of fitted models such that, for each $m$, Assumption~\ref{ass:monotone_regressivity} holds on $S$ and
\[
\max_{i\in S}|r_i(f_m)-1|\to 0.
\]
Then
\[
\Psi_S^{\mathrm{surr}}(f_m)\to 0,
\qquad
\widehat{\operatorname{Cov}}_S(r(f_m),z)\to 0,
\qquad
\widehat{\operatorname{Cov}}_S(u(f_m),z)\to 0,
\]
and the assessor-facing vertical-equity metrics converge to their neutral values:
\[
\operatorname{PRD}(S)\to 1,
\qquad
\operatorname{PRB}(S)\to 0,
\qquad
\operatorname{VEI}(S)\to 0.
\]
\end{corollary}

\begin{proof}
Uniform convergence $r_i(f_m)\to 1$ on the finite sample $S$ implies $e_i(f_m)=r_i(f_m)-1\to 0$ uniformly, and therefore $\Psi_S^{\mathrm{surr}}(f_m)\to 0$ by \eqref{eq:surrogate_penalty}. The covariances with $z$ also converge to zero because they are continuous finite-sample functionals of $r_i(f_m)$ and $u_i(f_m)=\log r_i(f_m)$.

For the metrics, $\operatorname{PRD}(S)\to 1$ follows from \eqref{eq:prd_cov_identity}. Likewise, the PRB slope converges to zero because its response variable converges uniformly to zero, and $\operatorname{VEI}(S)\to 0$ because the medians in the top and bottom value groups both converge to $1$.
\end{proof}

\begin{remark}[Why the monotone assumption is needed]
\label{rem:cov_not_sufficient}
A small covariance by itself does not rule out more complex nonmonotone price effects, because positive and negative local trends can cancel in the aggregate. The role of Assumption~\ref{ass:monotone_regressivity} is precisely to align the covariance target with the assessor metrics. Accordingly, the covariance penalty proposed here should not be interpreted as a replacement for PRD, PRB, or VEI, but rather as a smooth training-time control for the first-order directional component of regressivity, whose reduction moves the standard vertical-equity diagnostics in the desired direction under the monotone root behavior described above.
\end{remark}




\section{Algorithmic Details for Covariance-Penalized Boosting}
\label{app:boosting_algorithms}

This appendix collects pseudocode for boosted-tree implementations of the squared covariance penalty introduced in Section~\ref{sec:covariance_penalizations}. Throughout, we use the manuscript's sign convention
\[
e_i = f(x_i)-y_i,
\qquad
c_i = y_i-\bar y,
\qquad
z = \widehat{\Cov}_n(e,y)=\frac{1}{n}\sum_{i=1}^n e_i c_i.
\]
Under this convention, the squared covariance penalty is
\[
\Psi_n^{\mathrm{cov}^2}(f)=\frac{1}{2}z^2.
\]

\subsection{Exact Newton-style tree boosting with squared covariance penalty}
\label{app:exact_newton_cov_boosting}

The following algorithm implements the exact leaf-wise Newton update discussed in Section~\ref{subsec:tree_boosting}. This is the most relevant pseudocode to retain if only one boosted-tree algorithm is included in the appendix.

\begin{algorithm}[!ht]
\caption{Exact Newton-style tree boosting with squared covariance penalty}
\label{alg:exact_newton_cov_boosting}
\begin{algorithmic}[1]
\Require Training data $\D=\{(x_i,y_i)\}_{i=1}^n$, penalty weight $\rho \ge 0$, leaf regularization $\lambda \ge 0$, tree complexity penalty $\gamma \ge 0$, learning rate $\nu \in (0,1]$, number of rounds $M$
\State Compute $\bar y \gets \frac{1}{n}\sum_{i=1}^n y_i$
\State Compute $c_i \gets y_i-\bar y$ for all $i=1,\dots,n$
\State Initialize $f_0(x) \gets \bar y$
\For{$m=1,\dots,M$}
    \State Compute current predictions $f_i \gets f_{m-1}(x_i)$ for all $i=1,\dots,n$
    \State Compute residuals $e_i \gets f_i-y_i$ for all $i=1,\dots,n$
    \State Compute covariance statistic
    \[
    z \gets \frac{1}{n}\sum_{i=1}^n e_i c_i
    \]
    \State Compute exact gradients
    \[
    g_i \gets \frac{e_i}{n} + \frac{\rho}{n} z c_i
    \qquad \forall i=1,\dots,n
    \]
    \State Search over candidate tree structures $T$
    \For{each candidate tree $T$ with leaves $I_1,\dots,I_J$}
        \For{$j=1,\dots,J$}
            \State $G_j \gets \sum_{i\in I_j} g_i$
            \State $n_j \gets |I_j|$
            \State $s_j \gets \sum_{i\in I_j} c_i$
        \EndFor
        \State Form
        \[
        D \gets \diag\!\left(\frac{n_1}{n}+\lambda,\dots,\frac{n_J}{n}+\lambda\right),
        \qquad
        \beta \gets \frac{\rho}{n^2},
        \qquad
        A_T \gets D + \beta ss^\top
        \]
        \State Compute exact tree gain
        \[
        \operatorname{Gain}(T)
        \gets
        \frac{1}{2}G^\top A_T^{-1}G - \gamma J
        \]
    \EndFor
    \State Select the best tree structure $T_m$
    \State For the selected tree, form $A_m \gets D + \beta ss^\top$
    \State Compute optimal leaf weights
    \[
    w^\star \gets -A_m^{-1}G
    \]
    \State Update
    \[
    f_m(\cdot) \gets f_{m-1}(\cdot) + \nu\,T_m(\cdot;w^\star)
    \]
\EndFor
\State \Return final predictor $f_M$
\end{algorithmic}
\end{algorithm}

\subsection{Optional first-order boosting variant}
\label{app:first_order_cov_boosting}

For completeness, we also record a first-order boosting variant based on fitting base learners to the negative gradient of the penalized objective. This algorithm is mainly pedagogical, and can be omitted if a shorter appendix is preferred.

\begin{algorithm}[!ht]
\caption{First-order gradient boosting with squared covariance penalty}
\label{alg:first_order_cov_boosting}
\begin{algorithmic}[1]
\Require Training data $\D=\{(x_i,y_i)\}_{i=1}^n$, penalty weight $\rho \ge 0$, learning rate $\nu \in (0,1]$, number of rounds $M$
\State Compute $\bar y \gets \frac{1}{n}\sum_{i=1}^n y_i$
\State Compute $c_i \gets y_i-\bar y$ for all $i=1,\dots,n$
\State Initialize $f_0(x) \gets \bar y$
\For{$m=1,\dots,M$}
    \State Compute current predictions $f_i \gets f_{m-1}(x_i)$ for all $i=1,\dots,n$
    \State Compute residuals $e_i \gets f_i-y_i$ for all $i=1,\dots,n$
    \State Compute covariance statistic
    \[
    z \gets \frac{1}{n}\sum_{i=1}^n e_i c_i
    \]
    \State Form negative-gradient targets
    \[
    u_i \gets -\bigl(e_i + \rho z c_i\bigr)
    \qquad \forall i=1,\dots,n
    \]
    \State Fit a base learner structure $T_m$ to $\{(x_i,u_i)\}_{i=1}^n$
    \If{$T_m$ is a generic weak learner with output vector $h \in \R^n$}
        \State Compute
        \[
        q \gets \frac{1}{n}c^\top h
        \]
        \State Set exact line-search step
        \[
        \gamma_m
        \gets
        -
        \frac{\frac{1}{n}e^\top h + \rho z q}
        {\frac{1}{n}\|h\|_2^2 + \rho q^2}
        \]
        \State Update
        \[
        f_m(\cdot) \gets f_{m-1}(\cdot) + \nu\,\gamma_m h(\cdot)
        \]
    \Else
        \State Let $I_1,\dots,I_J$ be the leaves of $T_m$
        \For{$j=1,\dots,J$}
            \State $n_j \gets |I_j|$
            \State $b_j \gets \sum_{i\in I_j} e_i$
            \State $s_j \gets \sum_{i\in I_j} c_i$
            \State $q_j \gets s_j/n$
        \EndFor
        \State Form
        \[
        D \gets \diag\!\left(\frac{n_1}{n},\dots,\frac{n_J}{n}\right),
        \qquad
        q \gets (q_1,\dots,q_J)^\top,
        \qquad
        \eta \gets \frac{1}{n}b + \rho z q
        \]
        \State Solve the coupled leaf system
        \[
        (D+\rho qq^\top)w = -\eta
        \]
        \State Update
        \[
        f_m(\cdot) \gets f_{m-1}(\cdot) + \nu\,T_m(\cdot;w)
        \]
    \EndIf
\EndFor
\State \Return final predictor $f_M$
\end{algorithmic}
\end{algorithm}


% \newpage
% \section{Bulletpoints}

% \begin{itemize}
%     \item \textbf{Why this problem?}
%     \begin{enumerate}
%         \item Mass Appraisal with ML has been recently advancing. Examples are the CCAO, 
%         \item Mistakes on it can generate highly unfair situations in terms of taxation; specially if regressive: low-valued houses are overcharged and high-values houses are undercharged (e.g.
%     \end{enumerate}

    
%     \item \textbf{Why this approach?}
% \end{itemize}

% ==========================================
% Table 1: Validation Split
% ==========================================
\begin{table}[!htbp]
\centering
\renewcommand{\arraystretch}{1.5} % Extra stretch for stacked cell contents
\caption{Validation Results: Metrics in bold show improvement over LGBMRegressor. Colored percentages in parentheses under other models indicate percentual change (improvement / worsening) compared to LGBMRegressor.}
\label{tab:validation_metrics}
\resizebox{\textwidth}{!}{%
\begin{tabular}{ll|rr|rr|rrr}
\toprule
Model & Config ID & R2 & MAE & COD & COV (IAAO) & VEI & PRD & PRB \\
\midrule
LinearRegression & fd63507d & 0.668 & \$92,314 & 32.716 & 0.490 & 8.425\% & 1.026 & 0.059 \\
LGBMRegressor & 0c59889c & 0.889 & \$58,430 & 23.283 & 0.377 & -19.961\% & 1.066 & -0.069 \\
\midrule
SINGLE\_OPT\_REQ\_ROBUST & - & 
\begin{tabular}{@{}c@{}} 0.881 \\ \small\color{worsRed}{(-0.90\%)} \end{tabular} & 
\begin{tabular}{@{}c@{}} \$63,918 \\ \small\color{worsRed}{(-9.39\%)} \end{tabular} & 
\begin{tabular}{@{}c@{}} 24.465 \\ \small\color{worsRed}{(-5.08\%)} \end{tabular} & 
\begin{tabular}{@{}c@{}} \textbf{0.361} \\ \small\color{impGreen}{(+4.24\%)} \end{tabular} & 
\begin{tabular}{@{}c@{}} \textbf{2.538\%} \\ \small\color{impGreen}{(+87.29\%)} \end{tabular} & 
\begin{tabular}{@{}c@{}} \textbf{1.009} \\ \small\color{impGreen}{(+86.36\%)} \end{tabular} & 
\begin{tabular}{@{}c@{}} \textbf{0.026} \\ \small\color{impGreen}{(+62.32\%)} \end{tabular} \\
SINGLE\_OPT\_UTOPIA\_ROBUST & - & 
\begin{tabular}{@{}c@{}} 0.881 \\ \small\color{worsRed}{(-0.90\%)} \end{tabular} & 
\begin{tabular}{@{}c@{}} \$63,918 \\ \small\color{worsRed}{(-9.39\%)} \end{tabular} & 
\begin{tabular}{@{}c@{}} 24.465 \\ \small\color{worsRed}{(-5.08\%)} \end{tabular} & 
\begin{tabular}{@{}c@{}} \textbf{0.361} \\ \small\color{impGreen}{(+4.24\%)} \end{tabular} & 
\begin{tabular}{@{}c@{}} \textbf{2.538\%} \\ \small\color{impGreen}{(+87.29\%)} \end{tabular} & 
\begin{tabular}{@{}c@{}} \textbf{1.009} \\ \small\color{impGreen}{(+86.36\%)} \end{tabular} & 
\begin{tabular}{@{}c@{}} \textbf{0.026} \\ \small\color{impGreen}{(+62.32\%)} \end{tabular} \\
\bottomrule
\end{tabular}%
}
\end{table}

\vspace{3em}

% ==========================================
% Table 2: Test Split
% ==========================================
\begin{table}[!htbp]
\centering
\renewcommand{\arraystretch}{1.5} % Extra stretch for stacked cell contents
\caption{Test Results: Metrics in bold show improvement over LGBMRegressor. Colored percentages in parentheses under other models indicate percentual change (improvement / worsening) compared to LGBMRegressor.}
\label{tab:test_metrics}
\resizebox{\textwidth}{!}{%
\begin{tabular}{ll|rr|rr|rrr}
\toprule
Model & Config ID & R2 & MAE & COD & COV (IAAO) & VEI & PRD & PRB \\
\midrule
LinearRegression & fd63507d & 0.762 & \$86,244 & 25.013 & 0.369 & -13.527\% & 1.039 & -0.011 \\
LGBMRegressor & 0c59889c & 0.882 & \$72,672 & 23.437 & 0.359 & -37.257\% & 1.088 & -0.120 \\
\midrule
STACKING\_OPT\_AVG\_ROBUST & - & 
\begin{tabular}{@{}c@{}} \textbf{0.883} \\ \small\color{impGreen}{(+0.11\%)} \end{tabular} & 
\begin{tabular}{@{}c@{}} \$73,080 \\ \small\color{worsRed}{(-0.56\%)} \end{tabular} & 
\begin{tabular}{@{}c@{}} \textbf{22.823} \\ \small\color{impGreen}{(+2.62\%)} \end{tabular} & 
\begin{tabular}{@{}c@{}} \textbf{0.335} \\ \small\color{impGreen}{(+6.69\%)} \end{tabular} & 
\begin{tabular}{@{}c@{}} \textbf{-13.100\%} \\ \small\color{impGreen}{(+64.84\%)} \end{tabular} & 
\begin{tabular}{@{}c@{}} \textbf{1.040} \\ \small\color{impGreen}{(+54.55\%)} \end{tabular} & 
\begin{tabular}{@{}c@{}} \textbf{-0.020} \\ \small\color{impGreen}{(+83.33\%)} \end{tabular} \\
STACKING\_OPT\_UTOPIA\_ROBUST & - & 
\begin{tabular}{@{}c@{}} 0.879 \\ \small\color{worsRed}{(-0.34\%)} \end{tabular} & 
\begin{tabular}{@{}c@{}} \$76,258 \\ \small\color{worsRed}{(-4.94\%)} \end{tabular} & 
\begin{tabular}{@{}c@{}} 23.636 \\ \small\color{worsRed}{(-0.85\%)} \end{tabular} & 
\begin{tabular}{@{}c@{}} \textbf{0.345} \\ \small\color{impGreen}{(+3.90\%)} \end{tabular} & 
\begin{tabular}{@{}c@{}} \textbf{-3.157\%} \\ \small\color{impGreen}{(+91.53\%)} \end{tabular} & 
\begin{tabular}{@{}c@{}} \textbf{1.027} \\ \small\color{impGreen}{(+69.32\%)} \end{tabular} & 
\begin{tabular}{@{}c@{}} \textbf{0.014} \\ \small\color{impGreen}{(+88.33\%)} \end{tabular} \\
SINGLE\_OPT\_REQ\_ROBUST & - & 
\begin{tabular}{@{}c@{}} 0.878 \\ \small\color{worsRed}{(-0.45\%)} \end{tabular} & 
\begin{tabular}{@{}c@{}} \$77,163 \\ \small\color{worsRed}{(-6.18\%)} \end{tabular} & 
\begin{tabular}{@{}c@{}} 23.929 \\ \small\color{worsRed}{(-2.10\%)} \end{tabular} & 
\begin{tabular}{@{}c@{}} \textbf{0.348} \\ \small\color{impGreen}{(+3.06\%)} \end{tabular} & 
\begin{tabular}{@{}c@{}} \textbf{-2.237\%} \\ \small\color{impGreen}{(+94.00\%)} \end{tabular} & 
\begin{tabular}{@{}c@{}} \textbf{1.025} \\ \small\color{impGreen}{(+71.59\%)} \end{tabular} & 
\begin{tabular}{@{}c@{}} \textbf{0.019} \\ \small\color{impGreen}{(+84.17\%)} \end{tabular} \\
SINGLE\_OPT\_UTOPIA\_ROBUST & - & 
\begin{tabular}{@{}c@{}} 0.878 \\ \small\color{worsRed}{(-0.45\%)} \end{tabular} & 
\begin{tabular}{@{}c@{}} \$77,163 \\ \small\color{worsRed}{(-6.18\%)} \end{tabular} & 
\begin{tabular}{@{}c@{}} 23.929 \\ \small\color{worsRed}{(-2.10\%)} \end{tabular} & 
\begin{tabular}{@{}c@{}} \textbf{0.348} \\ \small\color{impGreen}{(+3.06\%)} \end{tabular} & 
\begin{tabular}{@{}c@{}} \textbf{-2.237\%} \\ \small\color{impGreen}{(+94.00\%)} \end{tabular} & 
\begin{tabular}{@{}c@{}} \textbf{1.025} \\ \small\color{impGreen}{(+71.59\%)} \end{tabular} & 
\begin{tabular}{@{}c@{}} \textbf{0.019} \\ \small\color{impGreen}{(+84.17\%)} \end{tabular} \\
\bottomrule
\end{tabular}%
}
\end{table}




\end{document}
